% !TEX program = xelatex
%==========================================================
% مبانی داده‌کاوی
% فایل اصلی کتاب — همه فصل‌ها از اینجا فراخوانی می‌شوند
%==========================================================
\documentclass[12pt,oneside]{memoir}

% ---------- بسته‌های پایه ----------
\usepackage{amsmath,amssymb,amsfonts}
\usepackage{graphicx}
\usepackage{booktabs}
\usepackage{longtable}
\usepackage{algorithm}
\usepackage{algpseudocode}
\usepackage{hyperref}
\usepackage{xcolor}
\usepackage{listings}     % برای کد پایتون در پیوست
\usepackage{caption}
\usepackage{tikz}
\usetikzlibrary{shapes.geometric,arrows.meta,positioning,calc}
% ---------- پشتیبانی فارسی (XeLaTeX) ----------
\usepackage{xepersian}
\settextfont{XB Niloofar}      % در صورت نبود فونت، با فونت فارسی نصب‌شده جایگزین کنید
\setlatintextfont{Times New Roman}

% ---------- مراجع ----------
\usepackage[backend=biber,style=numeric,sorting=none]{biblatex}
\addbibresource{bibliography/references.bib}

% ---------- اصطلاحات بیگانه ----------
% \dmterm{معادل فارسی}{اصطلاح انگلیسی}
% متن فارسی را در بدنه نشان می‌دهد و اصطلاح انگلیسی را به‌صورت زیرنویس
% و با جهت صحیح (چپ‌به‌راست) می‌آورد. برای اولین‌بار که یک اصطلاح در هر
% فصل معرفی می‌شود استفاده شود.
\newcommand{\dmterm}[2]{#1\footnote{\LR{#2}}}

% ---------- تنظیمات کد پایتون ----------
\lstdefinestyle{pythonstyle}{
    language=Python,
    basicstyle=\ttfamily\small %\LTRfootnote,
    keywordstyle=\color{blue},
    commentstyle=\color{gray},
    stringstyle=\color{teal},
    numbers=left,
    numberstyle=\tiny,
    breaklines=true,
    frame=single
}
\lstset{style=pythonstyle}

% ---------- اطلاعات کتاب ----------
\title{مبانی داده‌کاوی}
\author{نویسنده: محرم منصوری زاده}
\date{\today}

\begin{document}

\frontmatter
\maketitle
\chapter*{پیشگفتار}
\addcontentsline{toc}{chapter}{پیشگفتار}

این کتاب با هدف آموزش مبانی داده‌کاوی به دانشجویان سال‌های پایانی کارشناسی نوشته شده است.
محتوای آن بر پایهٔ دو مرجع استاندارد این حوزه تنظیم شده:

\begin{itemize}
    \item Han, J., Kamber, M., \& Pei, J. \textit{Data Mining: Concepts and Techniques}.
    \item Tan, P.-N., Steinbach, M., \& Kumar, V. \textit{Introduction to Data Mining}.
\end{itemize}

برخلاف مراجع اصلی که دانشنامه‌ای و جامع هستند، این کتاب برای یک درس ترمی طراحی شده و
تلاش شده مطالب به‌صورت فشرده و کاربردی ارائه شوند. در پایان هر فصل، پیاده‌سازی مثال‌ها
در زبان پایتون (پیوست الف) آمده است تا خواننده بتواند مفاهیم نظری را با کد عملی تطبیق دهد.

این کتاب به‌صورت متن‌باز روی گیت‌هاب منتشر می‌شود و از بازخورد، اصلاحیه و مشارکت خوانندگان استقبال می‌کند.
 \newpage
\tableofcontents

\mainmatter

% ---------- بخش اول: مبانی ----------
\part{مبانی}
\chapter{مقدمه‌ای بر داده‌کاوی}
\label{ch:intro}

\section*{اهداف این فصل}
\begin{itemize}
    \item درک تفاوت داده‌کاوی با آمار سنتی، پایگاه داده و یادگیری ماشین
    \item آشنایی با فرآیند کشف دانش از داده و جایگاه داده‌کاوی در آن
    \item شناخت انواع الگوهایی که می‌توان از داده استخراج کرد
    \item آشنایی با کاربردهای واقعی و چالش‌های اصلی داده‌کاوی
\end{itemize}

\section{داده‌کاوی چیست؟}

در دهه‌های اخیر، حجم داده‌های تولیدشده توسط سازمان‌ها، وب‌سایت‌ها، شبکه‌های اجتماعی،
حسگرها و سامانه‌های تراکنشی به شکلی بی‌سابقه رشد کرده است. بانک‌ها میلیون‌ها تراکنش
روزانه ثبت می‌کنند، فروشگاه‌های آنلاین رفتار خرید هر کاربر را ذخیره می‌کنند و شبکه‌های
اجتماعی روزانه میلیاردها پیام و تعامل تولید می‌کنند. مسئلهٔ اصلی این نیست که داده کم
است؛ برعکس، مسئله این است که در میان این حجم عظیم داده، دانش مفید و قابل‌اتکا چگونه
استخراج شود. این دقیقاً همان مسئله‌ای است که \dmterm{\textbf{داده‌کاوی}}{Data Mining}
به آن می‌پردازد.

داده‌کاوی را می‌توان این‌گونه تعریف کرد:

\begin{quote}
داده‌کاوی فرآیند کشف الگوها، روابط و دانش پنهان، غیربدیهی و بالقوه مفید از میان
حجم زیادی از داده است.
\end{quote}

سه ویژگی در این تعریف کلیدی هستند:

\begin{itemize}
    \item \textbf{پنهان بودن}: الگوهایی که به‌راحتی و با نگاه ساده به داده قابل مشاهده نیستند.
    \item \textbf{غیربدیهی بودن}: نتیجه نباید چیزی باشد که از قبل بدیهی یا شناخته‌شده است
          (مثلاً «مشتریانی که نان می‌خرند، معمولاً نان می‌خواهند» یک الگوی بی‌ارزش است).
    \item \textbf{مفید بودن}: الگوی کشف‌شده باید قابل استفاده در تصمیم‌گیری یا پیش‌بینی باشد.
\end{itemize}

\subsection{داده‌کاوی در برابر آمار، پایگاه داده و یادگیری ماشین}

سؤالی که معمولاً برای دانشجو پیش می‌آید این است: داده‌کاوی چه تفاوتی با آماری که
پیش از این خوانده‌اید، یا با پایگاه داده و یادگیری ماشین دارد؟ پاسخ این است که
داده‌کاوی یک حوزهٔ میان‌رشته‌ای است که از هر سه این حوزه‌ها بهره می‌برد:

\begin{itemize}
    \item از \textbf{آمار}، مبانی نظری برای اندازه‌گیری عدم قطعیت، آزمون فرضیه و
          برازش مدل را وام می‌گیرد.
    \item از \textbf{سیستم‌های پایگاه داده}، فناوری لازم برای ذخیره، بازیابی و
          مدیریت کارآمد حجم بزرگی از داده را می‌گیرد که اغلب در حافظه جا نمی‌شود.
    \item از \textbf{یادگیری ماشین} و هوش مصنوعی، الگوریتم‌هایی برای یادگیری خودکار
          الگو از داده به‌جای برنامه‌نویسی دستی قوانین را به کار می‌گیرد.
\end{itemize}

تفاوت اصلی داده‌کاوی با یادگیری ماشین صرف، در تمرکز آن است: یادگیری ماشین بیشتر
به الگوریتم و بهینه‌سازی مدل می‌پردازد، در حالی که داده‌کاوی کل فرآیند — از
داده خام تا دانش قابل استفاده — را در نظر می‌گیرد. جدول \ref{tab:dm-comparison}
این مقایسه را به‌طور خلاصه نشان می‌دهد.

\begin{table}[htbp]
\centering
\caption{مقایسهٔ داده‌کاوی با حوزه‌های نزدیک به آن}
\label{tab:dm-comparison}
\begin{tabular}{@{}p{2.6cm}p{5.8cm}p{4.6cm}@{}}
\toprule
\textbf{حوزه} & \textbf{تمرکز اصلی} & \textbf{سهم آن در داده‌کاوی} \\
\midrule
آمار & اندازه‌گیری عدم قطعیت، آزمون فرضیه، برازش مدل & مبانی نظری استنباط از داده \\
\addlinespace
پایگاه داده & ذخیره، بازیابی و مدیریت کارآمد داده حجیم & زیرساخت دسترسی به داده \\
\addlinespace
یادگیری ماشین & الگوریتم‌های یادگیری خودکار الگو از داده & موتور کشف الگو \\
\addlinespace
داده‌کاوی & کل فرآیند: از داده خام تا دانش قابل استفاده & ترکیب هر سه حوزهٔ بالا \\
\bottomrule
\end{tabular}
\end{table}

\section{فرآیند کشف دانش از داده}
\label{sec:kdd}

داده‌کاوی معمولاً به‌اشتباه به‌عنوان مترادف کل فرآیند تحلیل داده به کار می‌رود،
در حالی که در واقع تنها یکی از مراحل یک فرآیند بزرگ‌تر به نام
\dmterm{\textbf{کشف دانش از داده}}{Knowledge Discovery in Databases (KDD)} است.
شکل \ref{fig:kdd-flow} این فرآیند هفت‌مرحله‌ای را نشان می‌دهد.

\begin{figure}[htbp]
\centering
\begin{tikzpicture}[
    node distance=0.55cm and 0cm,
    box/.style={rectangle, rounded corners, draw=blue!60, thick,
                fill=blue!8, minimum width=8.5cm, minimum height=0.9cm,
                align=center, font=\small},
    arr/.style={-{Latex[length=2mm]}, thick, draw=blue!60}
]
\node[box] (n1) {گام ۱: پاک‌سازی داده};
\node[box, below=of n1] (n2) {گام ۲: یکپارچه‌سازی داده};
\node[box, below=of n2] (n3) {گام ۳: انتخاب داده};
\node[box, below=of n3] (n4) {گام ۴: تبدیل داده};
\node[box, below=of n4, fill=orange!15, draw=orange!70] (n5) {گام ۵: داده‌کاوی};
\node[box, below=of n5] (n6) {گام ۶: ارزیابی الگو};
\node[box, below=of n6, fill=green!10, draw=green!60!black] (n7) {گام ۷: ارائهٔ دانش};

\draw[arr] (n1) -- (n2);
\draw[arr] (n2) -- (n3);
\draw[arr] (n3) -- (n4);
\draw[arr] (n4) -- (n5);
\draw[arr] (n5) -- (n6);
\draw[arr] (n6) -- (n7);
\end{tikzpicture}
\caption{مراحل هفت‌گانهٔ فرآیند کشف دانش از داده}
\label{fig:kdd-flow}
\end{figure}

جزئیات هر مرحله به شرح زیر است:

\begin{enumerate}
    \item \dmterm{\textbf{پاک‌سازی داده}}{Data Cleaning}: حذف نویز و داده‌های ناسازگار.
    \item \dmterm{\textbf{یکپارچه‌سازی داده}}{Data Integration}: ترکیب داده از چند منبع مختلف.
    \item \dmterm{\textbf{انتخاب داده}}{Data Selection}: بازیابی داده‌های مرتبط با هدف تحلیل.
    \item \dmterm{\textbf{تبدیل داده}}{Data Transformation}: تبدیل داده به قالبی مناسب برای کاوش
          (مثلاً از طریق تجمیع یا نرمال‌سازی).
    \item \textbf{داده‌کاوی}: به‌کارگیری روش‌های هوشمند برای استخراج الگو.
    \item \dmterm{\textbf{ارزیابی الگو}}{Pattern Evaluation}: شناسایی الگوهای واقعاً جالب بر اساس
          معیارهایی مانند \dmterm{جالب بودن}{Interestingness}.
    \item \dmterm{\textbf{ارائهٔ دانش}}{Knowledge Presentation}: نمایش دانش کشف‌شده به کاربر با
          استفاده از تکنیک‌های تجسم داده و گزارش‌گیری.
\end{enumerate}

نکتهٔ مهم این است که در عمل، این مراحل خطی نیستند؛ فرآیند کشف دانش تکرارشونده است و
معمولاً پس از ارزیابی نتایج اولیه، لازم است به مراحل پیش‌تر (مثلاً انتخاب داده
یا تبدیل داده) بازگشت. در بسیاری از منابع صنعتی، نسخهٔ ساده‌شده‌تری از این
چرخه با نام \dmterm{\LR{CRISP-DM}}{Cross-Industry Standard Process for Data Mining}
نیز رایج است که شش مرحله دارد و به‌جای یک خط مستقیم، به‌صورت یک \textbf{چرخه}
دیده می‌شود (شکل \ref{fig:crispdm-cycle})، زیرا معمولاً پس از استقرار یک مدل در
عمل، پروژه دوباره از مرحلهٔ درک کسب‌وکار آغاز می‌شود.

\begin{figure}[htbp]
\centering
\begin{tikzpicture}[
    every node/.style={font=\scriptsize},
    circnode/.style={rectangle, rounded corners, draw=purple!70, thick,
                      fill=purple!8, minimum width=2.7cm, minimum height=0.85cm, align=center}
]
\def\n{6}
\def\R{3.3cm}
\node[circnode] (c0) at (90:\R)    {درک کسب‌وکار};
\node[circnode] (c1) at (30:\R)    {درک داده};
\node[circnode] (c2) at (-30:\R)   {آماده‌سازی داده};
\node[circnode] (c3) at (-90:\R)   {مدل‌سازی};
\node[circnode] (c4) at (-150:\R)  {ارزیابی};
\node[circnode] (c5) at (150:\R)   {استقرار};

\draw[-{Latex[length=2mm]}, thick, purple!70] (c0) -- (c1);
\draw[-{Latex[length=2mm]}, thick, purple!70] (c1) -- (c2);
\draw[-{Latex[length=2mm]}, thick, purple!70] (c2) -- (c3);
\draw[-{Latex[length=2mm]}, thick, purple!70] (c3) -- (c4);
\draw[-{Latex[length=2mm]}, thick, purple!70] (c4) -- (c5);
\draw[-{Latex[length=2mm]}, thick, purple!70] (c5) -- (c0);
\end{tikzpicture}
\caption{چرخهٔ \LR{CRISP-DM} به‌عنوان نسخهٔ صنعتی و تکرارشوندهٔ فرآیند کشف دانش}
\label{fig:crispdm-cycle}
\end{figure}

\subsection{چرا پیش‌پردازش این‌قدر مهم است؟}

تجربهٔ عملی نشان می‌دهد که در اغلب پروژه‌های داده‌کاوی، بیش از ۶۰ تا ۷۰ درصد
زمان صرف مراحل پاک‌سازی، یکپارچه‌سازی و تبدیل داده می‌شود، نه اجرای الگوریتم
یادگیری. دلیل آن ساده است: کیفیت هر مدل داده‌کاوی مستقیماً به کیفیت داده‌ای
که به آن داده می‌شود بستگی دارد؛ اصلی که با عبارت معروف «زباله وارد شود،
زباله بیرون می‌آید»\footnote{\LR{Garbage In, Garbage Out}} شناخته می‌شود. جزئیات
این مرحله در فصل \ref{ch:data} بررسی خواهد شد.

\section{انواع الگوهای قابل کشف}

خروجی داده‌کاوی می‌تواند اشکال گوناگونی داشته باشد. به‌طور کلی، وظایف داده‌کاوی
به دو دستهٔ اصلی تقسیم می‌شوند:

\begin{itemize}
    \item \dmterm{\textbf{وظایف توصیفی}}{Descriptive}: هدف، شناخت ویژگی‌های کلی داده است،
          بدون آنکه لزوماً پیش‌بینی خاصی صورت گیرد. نمونه‌ها: خوشه‌بندی، قوانین
          انجمنی، تشخیص داده پرت.
    \item \dmterm{\textbf{وظایف پیش‌بینی‌کننده}}{Predictive}: هدف، استفاده از برخی
          ویژگی‌ها برای پیش‌بینی مقدار یک ویژگی دیگر است. نمونه‌ها: طبقه‌بندی،
          رگرسیون.
\end{itemize}

جدول \ref{tab:dm-tasks} خلاصه‌ای از وظایف اصلی داده‌کاوی و فصل مرتبط با هر یک
در این کتاب را نشان می‌دهد.

\begin{table}[htbp]
\centering
\caption{وظایف اصلی داده‌کاوی و فصل مرتبط با هر یک در این کتاب}
\label{tab:dm-tasks}
\begin{tabular}{@{}lll@{}}
\toprule
\textbf{وظیفه} & \textbf{نوع} & \textbf{فصل مرتبط} \\
\midrule
طبقه‌بندی (\LR{Classification}) & پیش‌بینی‌کننده & فصل \ref{ch:classification-basic}، \ref{ch:classification-advanced} \\
رگرسیون (\LR{Regression}) & پیش‌بینی‌کننده & فصل \ref{ch:regression} \\
قوانین انجمنی (\LR{Association Rules}) & توصیفی & فصل \ref{ch:association} \\
خوشه‌بندی (\LR{Clustering}) & توصیفی & فصل \ref{ch:clustering} \\
تشخیص داده پرت (\LR{Outlier Detection}) & توصیفی/پیش‌بینی‌کننده & فصل \ref{ch:outliers} \\
\bottomrule
\end{tabular}
\end{table}

\section{کاربردهای داده‌کاوی}

داده‌کاوی امروزه در تقریباً هر صنعتی کاربرد دارد. جدول \ref{tab:dm-applications}
چند نمونهٔ رایج آن را نشان می‌دهد.

\begin{table}[htbp]
\centering
\caption{نمونه‌هایی از کاربرد داده‌کاوی در صنایع مختلف}
\label{tab:dm-applications}
\begin{tabular}{@{}p{3.2cm}p{9.8cm}@{}}
\toprule
\textbf{حوزه} & \textbf{کاربرد نمونه} \\
\midrule
بانکداری و مالی & تشخیص تقلب در تراکنش‌های کارت اعتباری، \dmterm{امتیازدهی اعتباری}{Credit Scoring} مشتریان \\
\addlinespace
خرده‌فروشی & \dmterm{تحلیل سبد خرید}{Market Basket Analysis} برای چیدمان محصولات و پیشنهاد خرید \\
\addlinespace
مراقبت سلامت & پیش‌بینی ریسک بیماری بر اساس سوابق بیمار، کشف الگوهای درمانی مؤثر \\
\addlinespace
بازاریابی & \dmterm{بخش‌بندی مشتریان}{Customer Segmentation}، \dmterm{سامانه‌های توصیه‌گر}{Recommender Systems} \\
\addlinespace
امنیت & \dmterm{تشخیص نفوذ به شبکه}{Intrusion Detection} \\
\bottomrule
\end{tabular}
\end{table}

\section{چالش‌های داده‌کاوی}

با وجود پیشرفت‌های چشمگیر، داده‌کاوی همچنان با چالش‌های جدی روبه‌روست:

\begin{itemize}
    \item \textbf{مقیاس‌پذیری}: بسیاری از الگوریتم‌های کلاسیک برای داده‌های
          کوچک طراحی شده‌اند و روی داده‌های بسیار بزرگ\footnote{\LR{Big Data}}
          کارایی کافی ندارند.
    \item \textbf{ابعاد بالا}: با افزایش تعداد ویژگی‌ها، بسیاری از روش‌ها
          دچار \dmterm{«نفرین ابعاد»}{Curse of Dimensionality} می‌شوند.
    \item \textbf{ناهمگونی و پیچیدگی داده}: داده‌های واقعی اغلب ترکیبی از
          متن، تصویر، سری زمانی و داده رابطه‌ای هستند.
    \item \textbf{حریم خصوصی و اخلاق داده}: استخراج الگو از داده‌های شخصی
          مسئولیت‌های حقوقی و اخلاقی جدی به همراه دارد.
    \item \textbf{تفسیرپذیری}: مدل‌های پیچیده (مانند جنگل تصادفی یا شبکه‌های
          عصبی) اغلب به‌صورت جعبه سیاه عمل می‌کنند و توضیح تصمیم آن‌ها دشوار است.
\end{itemize}

\section*{پیاده‌سازی در پایتون}

برای آشنایی عملی، مجموعه دادهٔ کلاسیک \LR{Iris} را بارگذاری کرده و یک
تحلیل اکتشافی دادهٔ ساده\footnote{\LR{Exploratory Data Analysis}} روی آن
انجام می‌دهیم: رسم پراکندگی
دو ویژگی «طول کاسبرگ» و «عرض کاسبرگ» برای سه گونهٔ گل، پیش از هرگونه
مدل‌سازی.

\begin{lstlisting}
import matplotlib.pyplot as plt
from sklearn.datasets import load_iris

iris = load_iris()
X, y, names = iris.data, iris.target, iris.target_names

fig, ax = plt.subplots(figsize=(5.5, 4.2))
for i, name in enumerate(names):
    mask = y == i
    ax.scatter(X[mask, 0], X[mask, 1], label=name.capitalize())

ax.set_xlabel("Sepal length (cm)")
ax.set_ylabel("Sepal width (cm)")
ax.legend()
plt.show()
\end{lstlisting}

خروجی این کد در شکل \ref{fig:iris-scatter} نشان داده شده است. حتی پیش از
اجرای هر الگوریتم داده‌کاوی، همین نمودار ساده به‌خوبی نشان می‌دهد که گونهٔ
\LR{Setosa} (نقاط آبی) بر اساس همین دو ویژگی به‌خوبی از دو گونهٔ دیگر
قابل تفکیک است، در حالی که \LR{Versicolor} و \LR{Virginica} همپوشانی
بیشتری دارند — دقیقاً نوع بینشی که تحلیل اکتشافی داده پیش از داده‌کاوی
باید فراهم کند.

\begin{figure}[htbp]
\centering
\includegraphics[width=0.72\linewidth]{figures/ch01_iris_exploratory.pdf}
\caption{نمودار پراکندگی طول و عرض کاسبرگ برای سه گونهٔ مجموعه دادهٔ Iris}
\label{fig:iris-scatter}
\end{figure}

کد کامل‌تر این فصل، به‌همراه چند نمودار اکتشافی دیگر، در بخش «فصل ۱» پیوست
الف (\ref{app:python}) آمده است.

\section*{خلاصهٔ فصل}

در این فصل با تعریف داده‌کاوی، جایگاه آن در فرآیند بزرگ‌تر کشف دانش، دسته‌بندی
وظایف اصلی آن به دو گروه توصیفی و پیش‌بینی‌کننده، و همچنین کاربردها و
چالش‌های اصلی این حوزه آشنا شدیم. در فصل بعد، به یکی از مهم‌ترین و
زمان‌برترین بخش‌های هر پروژهٔ داده‌کاوی، یعنی آشنایی با داده و پیش‌پردازش
آن، خواهیم پرداخت.

\section*{تمرین‌ها}

\begin{enumerate}
    \item تفاوت میان یک الگوی «بدیهی» و یک الگوی «جالب» را با یک مثال از
          حوزهٔ دلخواه خودتان توضیح دهید.
    \item برای هر یک از مراحل هفت‌گانهٔ کشف دانش (شکل \ref{fig:kdd-flow})، یک
          مثال عملی از یک پروژهٔ فرضی (مثلاً تحلیل رفتار مشتریان یک
          فروشگاه آنلاین) بیاورید.
    \item چرا چرخهٔ \LR{CRISP-DM} (شکل \ref{fig:crispdm-cycle}) به‌صورت
          دایره‌ای ترسیم می‌شود و نه به‌صورت یک خط مستقیم مانند فرآیند
          کشف دانش؟
    \item دو مورد از وظایف داده‌کاوی (مثلاً طبقه‌بندی و خوشه‌بندی) را از
          نظر نوع (توصیفی/پیش‌بینی‌کننده) و نوع داده‌ای که نیاز دارند
          (برچسب‌دار یا بدون برچسب) با هم مقایسه کنید.
    \item کد بخش «پیاده‌سازی در پایتون» را اجرا کنید و به‌جای دو ویژگی
          «طول و عرض کاسبرگ»، «طول و عرض گلبرگ» را رسم کنید. آیا تفکیک
          سه گونه در این حالت بهتر است یا بدتر؟ چرا؟
\end{enumerate}

\chapter{آشنایی با داده و پیش‌پردازش}
\label{ch:data}

\section*{اهداف این فصل}
\begin{itemize}
    \item شناخت انواع ویژگی‌ها و مقیاس‌های اندازه‌گیری داده
    \item درک مفهوم کیفیت داده و منابع اصلی مشکلات آن
    \item تسلط بر روش‌های پاک‌سازی، یکپارچه‌سازی، تبدیل و کاهش داده
    \item آشنایی با معیارهای شباهت و فاصله میان نمونه‌ها
    \item پیاده‌سازی عملی هر مرحله با پایتون و مشاهدهٔ اثر آن روی داده
\end{itemize}

\section{انواع ویژگی‌ها}
\label{sec:attribute-types}

هر مجموعه داده از تعدادی \dmterm{\textbf{نمونه}}{Instance / Object / Record}
تشکیل شده که هر یک با مجموعه‌ای از \dmterm{\textbf{ویژگی‌ها}}{Attributes / Features}
توصیف می‌شود. شناخت دقیق نوع هر ویژگی، پیش‌نیاز انتخاب روش پیش‌پردازش و الگوریتم
مناسب است. ویژگی‌ها را می‌توان از دو منظر دسته‌بندی کرد.

\subsection{دسته‌بندی کیفی در برابر کمّی}

\begin{itemize}
    \item \dmterm{\textbf{ویژگی‌های کیفی}}{Categorical / Qualitative}
    \begin{itemize}
        \item \dmterm{\textit{اسمی}}{Nominal}: مقادیر فقط برچسب‌اند و ترتیب یا فاصلهٔ
              معناداری ندارند. مثال: رنگ چشم، شهر محل سکونت.
        \item \dmterm{\textit{ترتیبی}}{Ordinal}: مقادیر ترتیب معناداری دارند، اما فاصلهٔ
              میان آن‌ها مشخص یا یکسان نیست. مثال: سطح تحصیلات (دیپلم، کارشناسی،
              کارشناسی‌ارشد)، رضایت مشتری (کم، متوسط، زیاد).
    \end{itemize}
    \item \dmterm{\textbf{ویژگی‌های کمّی}}{Numeric / Quantitative}
    \begin{itemize}
        \item \dmterm{\textit{بازه‌ای}}{Interval}: فاصله میان مقادیر معنادار است، اما
              صفر مطلق ندارد. مثال: دما بر حسب سلسیوس.
        \item \dmterm{\textit{نسبتی}}{Ratio}: علاوه بر فاصلهٔ معنادار، صفر مطلق نیز دارد،
              بنابراین نسبت دو مقدار هم معنادار است. مثال: سن، درآمد، وزن.
    \end{itemize}
\end{itemize}

هر نوع ویژگی مجموعهٔ مشخصی از عملیات ریاضی معنادار را مجاز می‌کند؛ برای مثال
میانگین‌گرفتن از «رنگ چشم» بی‌معناست، اما برای «سن» کاملاً معنادار است. جدول
\ref{tab:attr-types} این ارتباط را خلاصه می‌کند.

\begin{table}[htbp]
\centering
\caption{انواع ویژگی‌ها و عملیات معنادار برای هر نوع}
\label{tab:attr-types}
\begin{tabular}{@{}p{2.0cm}p{3.9cm}p{3.4cm}p{2.7cm}@{}}
\toprule
\textbf{نوع} & \textbf{ویژگی کلیدی} & \textbf{مثال} & \textbf{عملیات مجاز} \\
\midrule
اسمی & فقط برچسب، بدون ترتیب & رنگ چشم، کد شهر & برابری، مد \\
\addlinespace
ترتیبی & دارای ترتیب، بدون فاصلهٔ معنادار & سطح تحصیلات، رتبه & مقایسه، میانه \\
\addlinespace
بازه‌ای & فاصله معنادار، بدون صفر مطلق & دمای سلسیوس، سال تقویمی & جمع، تفریق، میانگین \\
\addlinespace
نسبتی & فاصله معنادار و صفر مطلق & سن، درآمد، وزن & همهٔ عملیات، نسبت \\
\bottomrule
\end{tabular}
\end{table}

\subsection{دسته‌بندی گسسته در برابر پیوسته}

از منظری دیگر، ویژگی‌ها به \textbf{گسسته} (تعداد محدود یا شمارش‌پذیر از
مقادیر، مانند تعداد فرزندان) و \textbf{پیوسته} (بی‌نهایت مقدار ممکن در یک
بازه، مانند قد یا وزن) تقسیم می‌شوند. تشخیص درست این نوع، در انتخاب روش
گسسته‌سازی (بخش \ref{subsec:discretization}) و همچنین نوع نمودار مناسب برای
تحلیل اکتشافی داده اهمیت دارد.

\section{کیفیت داده}
\label{sec:data-quality}

پیش از هر کاوشی، باید کیفیت داده ارزیابی شود. مهم‌ترین مشکلات رایج در داده‌های
واقعی عبارت‌اند از:

\begin{itemize}
    \item \dmterm{\textbf{مقادیر گمشده}}{Missing Values}: به دلایلی مانند خطای ثبت،
          عدم پاسخ‌گویی کاربر یا خرابی حسگر.
    \item \dmterm{\textbf{نویز}}{Noise}: خطای تصادفی یا انحراف در مقدار اندازه‌گیری‌شده.
    \item \dmterm{\textbf{داده‌های ناسازگار}}{Inconsistent Data}: تناقض میان رکوردها،
          مثلاً دو مقدار متفاوت برای سن یک فرد در دو جدول مختلف.
    \item \dmterm{\textbf{داده‌های تکراری}}{Duplicate Data}: ثبت چندبارهٔ یک رکورد.
    \item \dmterm{\textbf{داده‌های پرت}}{Outliers}: مقادیری که به‌طور قابل‌توجهی با
          بقیهٔ داده تفاوت دارند (جزئیات بیشتر در فصل \ref{ch:outliers}).
\end{itemize}

نکتهٔ مهم این است که نویز و داده پرت یکسان نیستند: نویز معمولاً خطای تصادفی
و بی‌معنی است که باید حذف یا اصلاح شود، اما داده پرت می‌تواند خودِ هدف تحلیل
باشد (مثلاً تراکنش متقلبانه).

اولین گام عملی در ارزیابی کیفیت، \textbf{دیدن} الگوی مقادیر گمشده است. شکل
\ref{fig:missing-values} نمونه‌ای از این نمودار را نشان می‌دهد: هر سطر یک رکورد،
هر ستون یک ویژگی، و خانه‌های سفید نشان‌دهندهٔ مقدار گمشده‌اند. در این مثال،
ستون‌های «درآمد» و «امتیاز رضایت» بیشترین گمشدگی را دارند و ستون «شهر» تقریباً
کامل است؛ چنین الگویی به تحلیلگر می‌گوید تلاش پاک‌سازی را روی کدام ستون‌ها
متمرکز کند.

\begin{figure}[htbp]
\centering
\includegraphics[width=0.62\linewidth]{figures/ch02_missing_values.pdf}
\caption{نمایش الگوی مقادیر گمشده در یک مجموعه داده (سفید = گمشده)}
\label{fig:missing-values}
\end{figure}

\section{پیش‌پردازش داده}
\label{sec:preprocessing}

پیش‌پردازش داده مجموعه‌ای از گام‌های پشت‌سرهم است که داده خام را به داده‌ای
مناسب برای الگوریتم‌های داده‌کاوی تبدیل می‌کند. شکل \ref{fig:preprocessing-pipeline}
ترتیب متداول این گام‌ها را نشان می‌دهد؛ هر یک در ادامه به‌تفصیل بررسی می‌شود.

\begin{figure}[htbp]
\centering
\begin{tikzpicture}[
    node distance=0.28cm,
    st/.style={rectangle, rounded corners, draw=blue!60, thick, fill=blue!8,
               minimum width=2.0cm, minimum height=0.95cm, text width=1.85cm,
               align=center, font=\scriptsize},
    endp/.style={rectangle, rounded corners, draw=gray!70, thick, fill=gray!12,
               minimum width=2.0cm, minimum height=0.95cm, text width=1.85cm,
               align=center, font=\scriptsize},
    note/.style={text width=2.0cm, align=center, font=\tiny, text=gray!70!black},
    arr/.style={-{Latex[length=2mm]}, thick, draw=blue!60}
]
\node[endp] (raw) {داده خام};
\node[st, left=of raw] (clean) {پاک‌سازی};
\node[st, left=of clean] (integ) {یکپارچه‌سازی};
\node[st, left=of integ] (trans) {تبدیل};
\node[st, left=of trans] (red) {کاهش};
\node[endp, left=of red, fill=green!10, draw=green!60!black] (ready) {داده آماده};

\node[note, below=0.15cm of clean] {حذف نویز و مقدار گمشده};
\node[note, below=0.15cm of integ] {ترکیب چند منبع داده};
\node[note, below=0.15cm of trans] {نرمال‌سازی و گسسته‌سازی};
\node[note, below=0.15cm of red] {نمونه‌برداری و \LR{PCA}};

\draw[arr] (raw) -- (clean);
\draw[arr] (clean) -- (integ);
\draw[arr] (integ) -- (trans);
\draw[arr] (trans) -- (red);
\draw[arr] (red) -- (ready);
\end{tikzpicture}
\caption{گام‌های متداول پیش‌پردازش داده}
\label{fig:preprocessing-pipeline}
\end{figure}

\subsection{پاک‌سازی داده}

رایج‌ترین راهکارها برای مقادیر گمشده در جدول \ref{tab:missing-strategies}
آمده‌اند. انتخاب میان آن‌ها به مقدار گمشدگی، علت آن و اهمیت ستون بستگی دارد.

\begin{table}[htbp]
\centering
\caption{راهکارهای مواجهه با مقادیر گمشده}
\label{tab:missing-strategies}
\begin{tabular}{@{}p{2.9cm}p{4.5cm}p{4.3cm}@{}}
\toprule
\textbf{راهکار} & \textbf{مزیت} & \textbf{عیب} \\
\midrule
حذف رکورد & ساده و سریع & اتلاف داده؛ در گمشدگی زیاد خطرناک \\
\addlinespace
جایگذاری با میانگین/میانه/مد & حفظ تعداد نمونه‌ها؛ پیاده‌سازی آسان & کاهش واریانس و نادیده‌گرفتن روابط میان ویژگی‌ها \\
\addlinespace
جایگذاری با مدل پیش‌بینی & استفاده از سایر ویژگی‌ها؛ دقت بالاتر & پیچیده‌تر و مستعد نشت اطلاعات \\
\bottomrule
\end{tabular}
\end{table}

برای هموارسازی نویز نیز روش‌هایی مانند \dmterm{\textbf{بین‌بندی}}{Binning}،
رگرسیون و تحلیل خوشه‌ای (برای شناسایی و حذف نقاط دورافتاده) به کار می‌روند.

نمونه‌کد زیر، شمارش مقادیر گمشده و جایگذاری آن‌ها با میانه را با
\texttt{pandas} و \texttt{scikit-learn} نشان می‌دهد:

\begin{lstlisting}
import pandas as pd
from sklearn.impute import SimpleImputer

df = pd.read_csv("customers.csv")

# 1) how many values are missing in each column?
print(df.isna().sum())

# 2) fill numeric gaps with the column median
num_cols = df.select_dtypes(include="number").columns
imputer = SimpleImputer(strategy="median")
df[num_cols] = imputer.fit_transform(df[num_cols])
\end{lstlisting}

\subsection{یکپارچه‌سازی داده}

وقتی داده از چند منبع مختلف (مثلاً چند پایگاه داده یا چند فایل) جمع‌آوری
می‌شود، باید به مسائلی مانند موارد زیر توجه کرد:

\begin{itemize}
    \item \dmterm{\textbf{تشخیص موجودیت یکسان}}{Entity Identification}: آیا
          \texttt{customer\_id} در یک جدول همان \texttt{cust\_no} در جدول
          دیگر است؟
    \item \dmterm{\textbf{افزونگی}}{Redundancy}: آیا یک ویژگی از ترکیب چند ویژگی
          دیگر قابل استخراج است؟ (مثلاً سود، از تفاضل درآمد و هزینه.)
    \item \textbf{تناقض مقیاس یا واحد}: مثلاً دما در یک منبع بر حسب سلسیوس و
          در منبع دیگر بر حسب فارنهایت ثبت شده باشد.
\end{itemize}

\subsection{تبدیل و نرمال‌سازی داده}

بسیاری از الگوریتم‌ها (به‌ویژه آن‌هایی که بر پایهٔ فاصله عمل می‌کنند، مانند
نزدیک‌ترین همسایه یا \LR{k-means}) به مقیاس ویژگی‌ها حساس‌اند. اگر یک ویژگی در
بازهٔ $[0, 1]$ و ویژگی دیگر در بازهٔ $[0, 100000]$ باشد، ویژگی دوم به‌طور
نادرست بر محاسبهٔ فاصله غالب می‌شود. دو روش رایج نرمال‌سازی عبارت‌اند از:

\paragraph{نرمال‌سازی \LR{Min-Max}}
مقدار هر ویژگی به بازهٔ $[0, 1]$ نگاشت می‌شود:
\[
x' = \frac{x - \min(x)}{\max(x) - \min(x)}
\]

\paragraph{استانداردسازی \LR{Z-score}}
مقدار هر ویژگی بر اساس میانگین ($\mu$) و انحراف معیار ($\sigma$) آن
تبدیل می‌شود:
\[
x' = \frac{x - \mu}{\sigma}
\]

استانداردسازی \LR{Z-score} معمولاً در حضور داده‌های پرت پایدارتر از
\LR{Min-Max} عمل می‌کند، زیرا \LR{Min-Max} به شدت به مقادیر حداقل و حداکثر
حساس است.

\paragraph{یک مثال عددی}
پنج مشاهده با مقادیر $10, 20, 30, 40, 100$ را در نظر بگیرید (مقدار ۱۰۰ یک
مقدار پرت نسبت به بقیه است). برای این داده $\min = 10$، $\max = 100$،
$\mu = 40$ و $\sigma \approx 31.62$ (انحراف معیار جامعه) است. نتیجهٔ دو
تبدیل در جدول \ref{tab:scaling-example} آمده است. دقت کنید که در \LR{Min-Max}
مقدار پرت باعث می‌شود چهار مقدار دیگر در بازهٔ باریک $[0, 0.33]$ فشرده شوند.

\begin{table}[htbp]
\centering
\caption{مقایسهٔ عددی نرمال‌سازی \LR{Min-Max} و \LR{Z-score} روی یک نمونهٔ کوچک}
\label{tab:scaling-example}
\begin{tabular}{@{}cccc@{}}
\toprule
\textbf{ردیف} & $x$ & \textbf{\LR{Min-Max}} & \textbf{\LR{Z-score}} \\
\midrule
۱ & 10  & 0.000 & $-0.949$ \\
۲ & 20  & 0.111 & $-0.632$ \\
۳ & 30  & 0.222 & $-0.316$ \\
۴ & 40  & 0.333 & \phantom{$-$}0.000 \\
۵ & 100 & 1.000 & \phantom{$-$}1.897 \\
\bottomrule
\end{tabular}
\end{table}

اهمیت مقیاس‌بندی را وقتی بهتر می‌بینیم که چند ویژگی با مقیاس‌های بسیار
متفاوت داشته باشیم. شکل \ref{fig:scaling-effect} دو خوشهٔ مصنوعی از افراد را
بر اساس «سن» (بازهٔ چند ده) و «درآمد» (بازهٔ ده‌ها هزار) نشان می‌دهد. در
نمودار سمت چپ (مقیاس خام)، فاصلهٔ اقلیدسی عملاً فقط به درآمد وابسته است و
سن هیچ نقشی ندارد؛ پس از هر دو نوع مقیاس‌بندی، دو ویژگی وزن قابل‌مقایسه‌ای
پیدا می‌کنند.

\begin{figure}[htbp]
\centering
\includegraphics[width=\linewidth]{figures/ch02_scaling_effect.pdf}
\caption{اثر مقیاس‌بندی بر دو ویژگی با مقیاس متفاوت: از چپ به راست، مقیاس خام،
\LR{Min-Max} و \LR{Z-score}}
\label{fig:scaling-effect}
\end{figure}

پیاده‌سازی این دو تبدیل با \texttt{scikit-learn} به‌صورت زیر است:

\begin{lstlisting}
import numpy as np
from sklearn.preprocessing import MinMaxScaler, StandardScaler

x = np.array([[10], [20], [30], [40], [100]], dtype=float)

x_minmax = MinMaxScaler().fit_transform(x)
x_zscore = StandardScaler().fit_transform(x)

print(x_minmax.ravel().round(3))  # [0.    0.111 0.222 0.333 1.   ]
print(x_zscore.ravel().round(3))  # [-0.949 -0.632 -0.316  0.     1.897]
\end{lstlisting}

\subsection{کاهش داده و کاهش ابعاد}

هدف کاهش داده، رسیدن به بازنمایی کوچک‌تر اما تقریباً برابر از نظر محتوای
اطلاعاتی است. دو رویکرد اصلی عبارت‌اند از:

\begin{itemize}
    \item \textbf{کاهش تعداد نمونه‌ها}: از طریق \dmterm{نمونه‌برداری}{Sampling} یا
          تجمیع داده.
    \item \textbf{کاهش تعداد ویژگی‌ها (کاهش ابعاد)}: از طریق
          \dmterm{انتخاب زیرمجموعه‌ای از ویژگی‌ها}{Feature Selection} یا
          استخراج ویژگی جدید.
\end{itemize}

رایج‌ترین روش استخراج ویژگی، \dmterm{\textbf{تحلیل مؤلفه‌های اصلی}}{Principal Component Analysis (PCA)}
است. ایدهٔ اصلی \LR{PCA} یافتن جهت‌هایی (مؤلفه‌های اصلی) در فضای ویژگی‌هاست که
بیشترین واریانس داده را حفظ می‌کنند؛ با تصویر کردن داده روی چند مؤلفهٔ اول،
می‌توان ابعاد را کاهش داد در حالی که بخش عمدهٔ اطلاعات داده حفظ می‌شود. توجه
شود که \LR{PCA} پیش از اجرا نیازمند استانداردسازی داده است، زیرا بر پایهٔ
واریانس عمل می‌کند.

مراحل کلی \LR{PCA} به‌صورت زیر است:

\begin{enumerate}
    \item استانداردسازی هر ویژگی (میانگین صفر و واریانس یک).
    \item محاسبهٔ ماتریس کوواریانس ویژگی‌ها.
    \item یافتن بردارهای ویژه و مقادیر ویژهٔ این ماتریس؛ بردار ویژه با بزرگ‌ترین
          مقدار ویژه، جهت بیشترین واریانس (مؤلفهٔ اول) است.
    \item انتخاب $k$ مؤلفهٔ اول و تصویر کردن داده روی آن‌ها.
\end{enumerate}

مجموعه دادهٔ \LR{Iris} چهار ویژگی دارد و نمی‌توان آن را مستقیماً روی صفحه
رسم کرد. شکل \ref{fig:pca-iris} نتیجهٔ کاهش آن به دو مؤلفهٔ اصلی را نشان می‌دهد.
دو مؤلفهٔ اول در مجموع حدود ۹۶ درصد واریانس داده را نگه می‌دارند
(۷۳ درصد و ۲۳ درصد)، و با وجود حذف دو بعد، سه گونه هنوز به‌خوبی قابل تشخیص‌اند.

\begin{figure}[htbp]
\centering
\includegraphics[width=0.72\linewidth]{figures/ch02_pca_iris.pdf}
\caption{تصویر مجموعه دادهٔ \LR{Iris} روی دو مؤلفهٔ اصلی (کاهش از چهار بعد به دو بعد)}
\label{fig:pca-iris}
\end{figure}

\begin{lstlisting}
from sklearn.datasets import load_iris
from sklearn.preprocessing import StandardScaler
from sklearn.decomposition import PCA

X = load_iris().data
X_scaled = StandardScaler().fit_transform(X)

pca = PCA(n_components=2)
X_2d = pca.fit_transform(X_scaled)

print(pca.explained_variance_ratio_)  # about [0.730 0.229]
\end{lstlisting}

\subsection{گسسته‌سازی و سلسله‌مراتب مفهومی}
\label{subsec:discretization}

برخی الگوریتم‌ها (مانند نسخه‌های پایهٔ درخت تصمیم یا بیز ساده) بهتر با
داده‌های کیفی کار می‌کنند. \textbf{گسسته‌سازی} فرآیند تبدیل یک ویژگی عددی
پیوسته به تعدادی بازهٔ گسسته (مانند تبدیل سن به گروه‌های «جوان، میان‌سال،
مسن») است. روش‌های رایج شامل بین‌بندی با عرض یکسان\footnote{\LR{Equal-width}}،
بین‌بندی با فراوانی یکسان\footnote{\LR{Equal-depth}} و روش‌های مبتنی بر
خوشه‌بندی هستند.

تفاوت این دو روش بین‌بندی در شکل \ref{fig:binning} دیده می‌شود. در
بین‌بندی با عرض یکسان، طول همهٔ بازه‌ها برابر است اما تعداد رکوردها در آن‌ها
یکسان نیست (وقتی داده چوله است، یک بازه ممکن است تقریباً خالی بماند). در
بین‌بندی با فراوانی یکسان، هر بازه تعداد برابری رکورد دارد اما عرض بازه‌ها
متفاوت است.

\begin{figure}[htbp]
\centering
\includegraphics[width=\linewidth]{figures/ch02_binning.pdf}
\caption{مقایسهٔ بین‌بندی با عرض یکسان (چپ) و بین‌بندی با فراوانی یکسان (راست) روی ویژگی سن}
\label{fig:binning}
\end{figure}

\begin{lstlisting}
import pandas as pd

ages = pd.Series([22, 25, 31, 34, 38, 41, 47, 52, 58, 66])

equal_width = pd.cut(ages, bins=4)    # same interval length
equal_depth = pd.qcut(ages, q=4)      # same number of records

print(equal_width.value_counts().sort_index())
print(equal_depth.value_counts().sort_index())
\end{lstlisting}

\dmterm{\textbf{سلسله‌مراتب مفهومی}}{Concept Hierarchy} نیز به تعریف سطوح مختلف
انتزاع برای یک ویژگی اشاره دارد (مثلاً شهر $\to$ استان $\to$ کشور)، که در
تحلیل‌های چندسطحی کاربرد دارد.

\section{معیارهای شباهت و فاصله}
\label{sec:similarity}

بسیاری از الگوریتم‌های داده‌کاوی (خوشه‌بندی، نزدیک‌ترین همسایه، تشخیص داده پرت)
بر پایهٔ سنجش شباهت یا فاصله میان دو نمونه عمل می‌کنند. رایج‌ترین معیارها
عبارت‌اند از:

\paragraph{فاصلهٔ اقلیدسی}
برای دو بردار $x$ و $y$ در فضای $n$-بعدی:
\[
d(x, y) = \sqrt{\sum_{i=1}^{n} (x_i - y_i)^2}
\]

\paragraph{فاصلهٔ منهتن}
\[
d(x, y) = \sum_{i=1}^{n} |x_i - y_i|
\]

تفاوت این دو فاصله را می‌توان روی صفحه دید. دو نقطهٔ $A=(0,0)$ و $B=(4,3)$ را
در نظر بگیرید. فاصلهٔ اقلیدسی طول پاره‌خط مستقیم میان آن‌هاست:
$\sqrt{4^2+3^2}=5$. فاصلهٔ منهتن مانند حرکت در خیابان‌های شبکه‌ای یک شهر است
که فقط در امتداد افقی و عمودی می‌توان رفت: $4+3=7$ (شکل \ref{fig:distances}).

\begin{figure}[htbp]
\centering
\begin{tikzpicture}[scale=0.95]
\draw[step=1, gray!30, very thin] (-0.4,-0.4) grid (5.4,4.4);
\draw[->, gray] (-0.4,0) -- (5.4,0);
\draw[->, gray] (0,-0.4) -- (0,4.4);
\coordinate (A) at (0,0);
\coordinate (B) at (4,3);
\draw[very thick, blue!70] (A) -- (B);
\draw[very thick, red!70, dashed]
    (0,0) -- (1,0) -- (1,1) -- (2,1) -- (2,2) -- (3,2) -- (3,3) -- (4,3);
\fill[black] (A) circle (2.5pt) node[below left] {$A$};
\fill[black] (B) circle (2.5pt) node[above right] {$B$};
\node[blue!70, font=\small] at (1.0,2.5) {$d_2 = 5$};
\node[red!70, font=\small] at (3.7,1.4) {$d_1 = 7$};

\draw[very thick, blue!70] (6.3,3.6) -- (7.0,3.6);
\node[right, font=\small] at (7.1,3.6) {فاصلهٔ اقلیدسی};
\draw[very thick, red!70, dashed] (6.3,2.9) -- (7.0,2.9);
\node[right, font=\small] at (7.1,2.9) {فاصلهٔ منهتن};
\end{tikzpicture}
\caption{مقایسهٔ فاصلهٔ اقلیدسی و فاصلهٔ منهتن میان دو نقطه در صفحه}
\label{fig:distances}
\end{figure}

\paragraph{فاصلهٔ مینکوفسکی}
تعمیمی از دو فاصلهٔ بالا با پارامتر $p$:
\[
d(x, y) = \left( \sum_{i=1}^{n} |x_i - y_i|^p \right)^{1/p}
\]
که برای $p=1$ به فاصلهٔ منهتن و برای $p=2$ به فاصلهٔ اقلیدسی تبدیل می‌شود.

\paragraph{شباهت کسینوسی}
برای داده‌های با ابعاد بالا و پراکنده (مانند بردارهای متنی)، شباهت کسینوسی
که زاویهٔ میان دو بردار را می‌سنجد، اغلب مناسب‌تر از فاصلهٔ اقلیدسی است:
\[
\text{sim}(x, y) = \frac{x \cdot y}{\|x\| \, \|y\|}
\]

برای ویژگی‌های اسمی یا دودویی نیز معیارهایی مانند
\dmterm{\textbf{ضریب تطابق ساده}}{Simple Matching Coefficient} و
\dmterm{\textbf{ضریب ژاکارد}}{Jaccard Coefficient} به کار می‌روند که در فصل
مربوط به خوشه‌بندی (فصل \ref{ch:clustering}) دوباره به آن‌ها پرداخته می‌شود.
جدول \ref{tab:distance-summary} چهار معیار اصلی این بخش را کنار هم قرار می‌دهد.

\begin{table}[htbp]
\centering
\caption{خلاصهٔ معیارهای فاصله و شباهت}
\label{tab:distance-summary}
\begin{tabular}{@{}p{2.6cm}p{3.9cm}p{5.6cm}@{}}
\toprule
\textbf{معیار} & \textbf{فرمول} & \textbf{کاربرد مناسب} \\
\midrule
اقلیدسی & $\sqrt{\sum_i (x_i-y_i)^2}$ & داده‌های عددی با ابعاد کم و مقیاس یکسان \\
\addlinespace
منهتن & $\sum_i |x_i-y_i|$ & داده‌هایی که مقادیر پرت دارند یا ابعاد آن‌ها مستقل‌اند \\
\addlinespace
مینکوفسکی & $\left(\sum_i |x_i-y_i|^p\right)^{1/p}$ & تعمیم دو معیار بالا؛ پارامتر $p$ قابل تنظیم \\
\addlinespace
کسینوسی & $\dfrac{x\cdot y}{\|x\|\,\|y\|}$ & داده‌های پراکنده و پرابعاد مانند متن \\
\bottomrule
\end{tabular}
\end{table}

محاسبهٔ این معیارها با \texttt{scipy} کوتاه است. برای دو نقطهٔ شکل
\ref{fig:distances}:

\begin{lstlisting}
from scipy.spatial import distance

A, B = [0, 0], [4, 3]

print(distance.euclidean(A, B))   # 5.0
print(distance.cityblock(A, B))   # 7   (Manhattan)
print(distance.minkowski(A, B, p=3))

# cosine *distance* = 1 - cosine similarity
print(1 - distance.cosine([1, 2, 3], [2, 4, 6]))  # 1.0
\end{lstlisting}

\section*{پیاده‌سازی در پایتون}

نمونه‌کدهای این فصل در متن هر بخش آمده‌اند. مجموعهٔ کامل آن‌ها به‌صورت یک
اسکریپت واحد (جایگذاری، مقیاس‌بندی، \LR{PCA}، بین‌بندی و محاسبهٔ فاصله) در بخش
«فصل ۲» پیوست الف (\ref{app:python}) گردآوری شده است.

\section*{خلاصهٔ فصل}

در این فصل با انواع ویژگی‌ها، منابع اصلی مشکلات کیفیت داده، و مراحل کلیدی
پیش‌پردازش شامل پاک‌سازی، یکپارچه‌سازی، تبدیل، کاهش داده و گسسته‌سازی آشنا
شدیم. همچنین رایج‌ترین معیارهای فاصله و شباهت را که در فصل‌های بعدی
(به‌ویژه طبقه‌بندی و خوشه‌بندی) بارها استفاده خواهند شد، مرور کردیم. در
فصل بعد، به اولین وظیفهٔ اصلی داده‌کاوی یعنی طبقه‌بندی می‌پردازیم.

\section*{تمرین‌ها}

\begin{enumerate}
    \item برای هر یک از ویژگی‌های «شمارهٔ ملی»، «دمای بدن»، «رتبهٔ مسابقه»
          و «تعداد فرزندان» نوع آن (اسمی/ترتیبی/بازه‌ای/نسبتی) را مشخص کنید.
    \item فرض کنید در یک مجموعه داده، ۱۵ درصد مقادیر ستون «درآمد» گمشده است.
          با توجه به جدول \ref{tab:missing-strategies}، دو راهکار پیشنهاد دهید و
          مزایا و معایب هرکدام را توضیح دهید.
    \item مقادیر $5, 15, 25, 35, 500$ را هم با \LR{Min-Max} و هم با \LR{Z-score}
          تبدیل کنید و نتیجه را با جدول \ref{tab:scaling-example} مقایسه کنید. کدام روش
          در برابر مقدار ۵۰۰ پایدارتر است؟
    \item فاصلهٔ اقلیدسی و فاصلهٔ منهتن را برای دو بردار
          $x = (2, 4, 6)$ و $y = (5, 1, 8)$ محاسبه کنید و پاسخ را با کد
          \texttt{scipy} بررسی کنید.
    \item چرا شباهت کسینوسی برای مقایسهٔ اسناد متنی معمولاً بهتر از فاصلهٔ
          اقلیدسی عمل می‌کند؟
    \item در شکل \ref{fig:pca-iris}، اگر فقط مؤلفهٔ اول را نگه می‌داشتیم، تفکیک کدام
          دو گونه دشوارتر می‌شد؟ چرا؟
    \item با \texttt{pd.cut} و \texttt{pd.qcut} یک ستون سن دلخواه را بین‌بندی کنید
          و تفاوت تعداد رکوردها در بازه‌ها را با شکل \ref{fig:binning} مقایسه کنید.
\end{enumerate}


% ---------- بخش دوم: یادگیری با نظارت ----------
\part{یادگیری با نظارت}
\chapter{طبقه‌بندی: مفاهیم پایه}
\label{ch:classification-basic}

\section*{اهداف این فصل}
\begin{itemize}
    \item درک مفهوم یادگیری با نظارت و مسئلهٔ طبقه‌بندی
    \item یادگیری نحوهٔ ساخت درخت تصمیم و معیارهای انتخاب ویژگی
    \item آشنایی با هرس درخت تصمیم برای جلوگیری از بیش‌برازش
    \item یادگیری طبقه‌بند بیز ساده
    \item آشنایی مقدماتی با ماتریس درهم‌ریختگی و معیار دقت
\end{itemize}

\section{مقدمه‌ای بر طبقه‌بندی}
\label{sec:classification-intro}

\dmterm{\textbf{طبقه‌بندی}}{Classification} یکی از رایج‌ترین وظایف
\dmterm{یادگیری با نظارت}{Supervised Learning} است. در این نوع یادگیری،
برخلاف خوشه‌بندی که در فصل \ref{ch:clustering} خواهیم دید، مجموعه‌ای از
نمونه‌های \dmterm{\textbf{برچسب‌دار}}{Labeled Data} در اختیار داریم؛ یعنی
برای هر نمونه، علاوه بر ویژگی‌های ورودی، مقدار کلاس هدف نیز مشخص است. هدف
طبقه‌بندی، یادگیری یک تابع $f: X \rightarrow Y$ است که بتواند برای
نمونه‌های جدید و دیده‌نشده، مقدار کلاس $Y$ را از روی ویژگی‌های $X$
پیش‌بینی کند.

مثال‌های رایج طبقه‌بندی عبارت‌اند از: تشخیص اسپم بودن یک ایمیل (بله/خیر)،
پیش‌بینی نکول وام یک متقاضی (پرداخت/نکول)، و تشخیص نوع گل بر اساس اندازهٔ
گلبرگ‌ها (مجموعه دادهٔ معروف \LR{Iris}).

فرآیند کلی ساخت یک مدل طبقه‌بندی دو مرحله دارد:

\begin{enumerate}
    \item \dmterm{\textbf{مرحلهٔ آموزش}}{Training}: مدل با استفاده از
          \dmterm{مجموعهٔ آموزش}{Training Set} که شامل نمونه‌های
          برچسب‌دار است، الگوی رابطهٔ میان ویژگی‌ها و کلاس را یاد می‌گیرد.
    \item \dmterm{\textbf{مرحلهٔ آزمون}}{Testing}: عملکرد مدل روی مجموعه‌ای
          جداگانه به نام \dmterm{مجموعهٔ آزمون}{Test Set} که در آموزش دیده
          نشده، سنجیده می‌شود تا از توانایی
          \dmterm{تعمیم}{Generalization} مدل مطمئن شویم.
\end{enumerate}

\section{درخت تصمیم}
\label{sec:decision-tree}

\dmterm{\textbf{درخت تصمیم}}{Decision Tree} یکی از قابل‌فهم‌ترین و
پرکاربردترین الگوریتم‌های طبقه‌بندی است. ساختار آن شبیه یک فلوچارت است:
هر گره داخلی یک آزمون روی یک ویژگی است، هر شاخه نتیجهٔ آن آزمون را نشان
می‌دهد و هر \dmterm{برگ}{Leaf} یک برچسب کلاس است. مزیت اصلی درخت تصمیم،
تفسیرپذیری بالای آن است؛ می‌توان مسیر از ریشه تا هر برگ را به‌صورت یک
قانون \LR{IF-THEN} ساده خواند.

سؤال اصلی در ساخت درخت تصمیم این است: در هر گره، کدام ویژگی برای تقسیم
داده انتخاب شود؟ پاسخ به این سؤال با معیارهای زیر داده می‌شود.

\subsection{معیارهای انتخاب ویژگی}
\label{subsec:split-criteria}

\paragraph{آنتروپی و بهرهٔ اطلاعاتی}
\dmterm{\textbf{بهرهٔ اطلاعاتی}}{Information Gain} معیاری است که در
الگوریتم \LR{ID3} استفاده می‌شود و بر پایهٔ مفهوم آنتروپی از نظریهٔ
اطلاعات ساخته شده است. برای یک مجموعه داده $D$ با $m$ کلاس، آنتروپی
به‌صورت زیر تعریف می‌شود:
\[
\text{Entropy}(D) = -\sum_{i=1}^{m} p_i \log_2(p_i)
\]
که در آن $p_i$ نسبت نمونه‌های کلاس $i$ در $D$ است. آنتروپی معیاری از
«ناخالصی» یا عدم قطعیت مجموعه داده است: هرچه کلاس‌ها یکنواخت‌تر توزیع
شده باشند، آنتروپی بیشتر است؛ اگر همهٔ نمونه‌ها به یک کلاس تعلق داشته
باشند، آنتروپی برابر صفر است.

پس از تقسیم داده بر اساس ویژگی $A$ به زیرمجموعه‌های $D_1, \ldots, D_v$،
بهرهٔ اطلاعاتی به‌صورت کاهش آنتروپی تعریف می‌شود:
\[
\text{Gain}(D, A) = \text{Entropy}(D) - \sum_{j=1}^{v} \frac{|D_j|}{|D|} \, \text{Entropy}(D_j)
\]
در هر گره، ویژگی‌ای انتخاب می‌شود که بیشترین بهرهٔ اطلاعاتی را دارد.

\paragraph{نسبت بهره}
یک ضعف بهرهٔ اطلاعاتی این است که به‌طور طبیعی به ویژگی‌هایی با تعداد
مقادیر زیاد (مانند شناسهٔ مشتری) گرایش پیدا می‌کند، حتی اگر آن ویژگی از
نظر پیش‌بینی بی‌ارزش باشد. الگوریتم \LR{C4.5} برای رفع این مشکل، معیار
\dmterm{\textbf{نسبت بهره}}{Gain Ratio} را معرفی کرد که بهرهٔ اطلاعاتی را
بر یک عامل نرمال‌ساز به نام
\dmterm{اطلاعات تقسیم}{Split Information} تقسیم می‌کند:
\[
\text{GainRatio}(D, A) = \frac{\text{Gain}(D, A)}{\text{SplitInfo}(D, A)}
\]

\paragraph{شاخص جینی}
\dmterm{\textbf{شاخص جینی}}{Gini Index} معیاری است که در الگوریتم
\LR{CART} استفاده می‌شود و احتمال طبقه‌بندی نادرست یک نمونهٔ تصادفی را
می‌سنجد:
\[
\text{Gini}(D) = 1 - \sum_{i=1}^{m} p_i^2
\]
مشابه آنتروپی، هرچه شاخص جینی کوچک‌تر باشد، مجموعه خالص‌تر است. در هر
گره، ویژگی‌ای انتخاب می‌شود که بیشترین کاهش شاخص جینی را ایجاد کند.

\subsection{هرس درخت تصمیم}
\label{subsec:pruning}

اگر اجازه داده شود درخت تا رسیدن به خلوص کامل هر برگ رشد کند، مدل به
جزئیات و نویز داده‌های آموزش بیش از حد وابسته می‌شود؛ پدیده‌ای که
\dmterm{\textbf{بیش‌برازش}}{Overfitting} نام دارد و باعث افت عملکرد روی
داده‌های جدید می‌شود. \dmterm{\textbf{هرس}}{Pruning} تکنیکی برای مقابله
با این مشکل است و به دو صورت انجام می‌شود:

\begin{itemize}
    \item \dmterm{\textbf{هرس پیشین}}{Pre-pruning}: رشد درخت زودتر متوقف
          می‌شود، مثلاً با تعیین حداکثر عمق درخت یا حداقل تعداد نمونه در
          هر برگ.
    \item \dmterm{\textbf{هرس پسین}}{Post-pruning}: ابتدا اجازه داده
          می‌شود درخت کامل رشد کند، سپس شاخه‌هایی که تأثیر چندانی بر دقت
          روی \dmterm{دادهٔ اعتبارسنجی}{Validation Set} ندارند، حذف
          می‌شوند.
\end{itemize}

در عمل، هرس پسین معمولاً نتایج بهتری می‌دهد، هرچند از نظر محاسباتی
پرهزینه‌تر است.

\section{طبقه‌بند بیز ساده}
\label{sec:naive-bayes}

\dmterm{\textbf{طبقه‌بند بیز ساده}}{Naive Bayes} روشی احتمالاتی برای
طبقه‌بندی است که بر پایهٔ قضیهٔ بیز عمل می‌کند:
\[
P(C \mid X) = \frac{P(X \mid C) \, P(C)}{P(X)}
\]
که در آن $C$ کلاس هدف و $X = (x_1, x_2, \ldots, x_n)$ بردار ویژگی‌های
یک نمونه است. هدف، یافتن کلاسی است که $P(C \mid X)$ را بیشینه می‌کند.

نام «ساده» (\LR{Naive}) به این دلیل است که این روش فرض می‌کند ویژگی‌ها،
مشروط بر کلاس، از یکدیگر \textbf{مستقل} هستند. با این فرض:
\[
P(X \mid C) = \prod_{i=1}^{n} P(x_i \mid C)
\]
با وجود ساده‌بودن (و اغلب نادرست‌بودن) این فرض استقلال در دنیای واقعی،
بیز ساده در عمل، به‌ویژه در طبقه‌بندی متن (مانند تشخیص اسپم)، عملکرد
شگفت‌آور خوبی دارد و از نظر محاسباتی بسیار کارآمد است، چون نیازی به
یادگیری روابط پیچیده میان ویژگی‌ها ندارد.

\paragraph{مسئلهٔ احتمال صفر و هموارسازی لاپلاس}
اگر یک مقدار ویژگی در دادهٔ آموزش هرگز با یک کلاس خاص هم‌رخداد نداشته
باشد، $P(x_i \mid C)$ برابر صفر محاسبه می‌شود و کل حاصل‌ضرب را صفر
می‌کند، حتی اگر سایر ویژگی‌ها به‌شدت بر آن کلاس دلالت کنند. برای رفع این
مشکل، از \dmterm{\textbf{هموارسازی لاپلاس}}{Laplace Smoothing} استفاده
می‌شود که یک مقدار کوچک به همهٔ شمارش‌ها اضافه می‌کند تا هیچ احتمالی
دقیقاً صفر نشود.

\section{ارزیابی مقدماتی}
\label{sec:basic-evaluation}

پس از آموزش یک مدل طبقه‌بندی، باید عملکرد آن را روی دادهٔ آزمون بسنجیم.
ابزار پایه برای این کار
\dmterm{\textbf{ماتریس درهم‌ریختگی}}{Confusion Matrix} است که برای یک
مسئلهٔ دوکلاسه به شکل زیر است:

\begin{center}
\begin{tabular}{@{}lcc@{}}
\toprule
 & \textbf{پیش‌بینی: مثبت} & \textbf{پیش‌بینی: منفی} \\
\midrule
\textbf{واقعی: مثبت} & \LR{True Positive (TP)} & \LR{False Negative (FN)} \\
\textbf{واقعی: منفی} & \LR{False Positive (FP)} & \LR{True Negative (TN)} \\
\bottomrule
\end{tabular}
\end{center}

ساده‌ترین معیار برخاسته از این ماتریس،
\dmterm{\textbf{دقت کلی}}{Accuracy} است:
\[
\text{Accuracy} = \frac{TP + TN}{TP + TN + FP + FN}
\]

هرچند دقت کلی معیاری شهودی است، اما در مسائلی که کلاس‌ها نامتوازن‌اند
(مثلاً تشخیص تقلب که تنها ۱ درصد تراکنش‌ها متقلبانه‌اند) می‌تواند
گمراه‌کننده باشد؛ مدلی که همیشه «غیرمتقلبانه» پیش‌بینی کند، دقت ۹۹ درصد
دارد اما کاملاً بی‌فایده است. معیارهای دقیق‌تر مانند \LR{Precision}،
\LR{Recall} و \LR{F1-Score}، که این مشکل را برطرف می‌کنند، در فصل
\ref{ch:classification-advanced} به‌طور کامل بررسی خواهند شد.

\section*{پیاده‌سازی در پایتون}

پیاده‌سازی درخت تصمیم (\texttt{DecisionTreeClassifier}) با معیارهای
\texttt{entropy} و \texttt{gini}، و همچنین طبقه‌بند \texttt{GaussianNB}،
به‌همراه رسم ماتریس درهم‌ریختگی، در بخش «فصل ۳» پیوست الف
(\ref{app:python}) آمده است.

\section*{خلاصهٔ فصل}

در این فصل با مفهوم یادگیری با نظارت و طبقه‌بندی آشنا شدیم. درخت تصمیم
را به‌عنوان یک مدل قابل‌تفسیر بررسی کردیم و سه معیار اصلی انتخاب ویژگی
(بهرهٔ اطلاعاتی، نسبت بهره و شاخص جینی) و همچنین اهمیت هرس برای جلوگیری
از بیش‌برازش را دیدیم. سپس طبقه‌بند بیز ساده را به‌عنوان یک روش
احتمالاتی و کارآمد معرفی کردیم. در پایان، با ماتریس درهم‌ریختگی و معیار
دقت، اولین ابزار ارزیابی مدل را آموختیم. در فصل بعد، به سراغ
الگوریتم‌های پیشرفته‌تر طبقه‌بندی (نزدیک‌ترین همسایه، ماشین بردار
پشتیبان، روش‌های ترکیبی) و معیارهای ارزیابی دقیق‌تر می‌رویم.

\section*{تمرین‌ها}

\begin{enumerate}
    \item آنتروپی یک مجموعه داده با ۱۰ نمونه که ۶ نمونه از کلاس مثبت و
          ۴ نمونه از کلاس منفی هستند را محاسبه کنید.
    \item چرا معیار بهرهٔ اطلاعاتی به‌تنهایی می‌تواند به ویژگی‌هایی با
          مقادیر بسیار زیاد (مانند شناسهٔ یکتای مشتری) گرایش پیدا کند؟
          نسبت بهره چگونه این مشکل را حل می‌کند؟
    \item تفاوت هرس پیشین و هرس پسین را از نظر هزینهٔ محاسباتی و کیفیت
          نتیجهٔ نهایی مقایسه کنید.
    \item فرض استقلال شرطی در بیز ساده را با یک مثال نقض کنید (یعنی
          مثالی از دو ویژگی که مستقل از هم نیستند) و توضیح دهید چرا با
          وجود این نقض، مدل همچنان می‌تواند عملکرد قابل‌قبولی داشته باشد.
    \item برای یک مسئلهٔ تشخیص بیماری نادر که تنها ۲ درصد جمعیت به آن
          مبتلا هستند، توضیح دهید چرا دقت کلی به‌تنهایی معیار مناسبی
          برای ارزیابی مدل نیست.
\end{enumerate}

\chapter{طبقه‌بندی پیشرفته و ارزیابی مدل}
\label{ch:classification-advanced}

\section*{اهداف این فصل}
\begin{itemize}
    \item آشنایی با k نزدیک‌ترین همسایه (k-NN) و ماشین بردار پشتیبان (SVM)
    \item یادگیری روش‌های ارکستراسیون: Bagging، Boosting، Random Forest
    \item تسلط بر معیارهای ارزیابی: Precision، Recall، F1، ROC/AUC
\end{itemize}

\section{k نزدیک‌ترین همسایه (k-NN)}
\section{ماشین بردار پشتیبان (SVM)}
\section{روش‌های ترکیبی (Ensemble)}
\subsection{Bagging و Random Forest}
\subsection{Boosting}

\section{ارزیابی مدل}
\subsection{اعتبارسنجی متقاطع (Cross-Validation)}
\subsection{Precision، Recall، F1-Score}
\subsection{منحنی ROC و AUC}
\subsection{مسئلهٔ عدم توازن کلاس‌ها (Class Imbalance)}

\section*{پیاده‌سازی در پایتون}
پیاده‌سازی k-NN، SVM و Random Forest با \texttt{scikit-learn} در پیوست الف (\ref{app:python}).

\section*{تمرین‌ها}

\chapter{مقدمه‌ای بر رگرسیون}
\label{ch:regression}

\section*{اهداف این فصل}
\begin{itemize}
    \item درک تفاوت رگرسیون با طبقه‌بندی
    \item یادگیری رگرسیون خطی ساده و چندگانه
    \item آشنایی با رگرسیون لجستیک به‌عنوان پل میان رگرسیون و طبقه‌بندی
    \item آشنایی با معیارهای ارزیابی مدل‌های رگرسیون
\end{itemize}

\section{رگرسیون در برابر طبقه‌بندی}

در دو فصل قبل، با \textbf{طبقه‌بندی} آشنا شدیم که در آن متغیر هدف
\textbf{کیفی} یا گسسته است (مثلاً «اسپم» یا «غیراسپم»).
\dmterm{\textbf{رگرسیون}}{Regression} نیز مانند طبقه‌بندی یک وظیفهٔ
یادگیری با نظارت است، با این تفاوت اساسی که متغیر هدف در آن
\textbf{عددی و پیوسته} است — مانند پیش‌بینی قیمت یک خانه، دمای فردا، یا
میزان فروش ماه آینده. بسیاری از مفاهیمی که در فصل‌های قبل دیدیم (تفکیک
آموزش/آزمون، اعتبارسنجی متقاطع، خطر بیش‌برازش) بدون تغییر برای رگرسیون
نیز صادق‌اند؛ آنچه تغییر می‌کند، نوع مدل و معیار ارزیابی است.

\section{رگرسیون خطی ساده و چندگانه}

\dmterm{\textbf{رگرسیون خطی ساده}}{Simple Linear Regression} رابطهٔ میان
یک متغیر مستقل $x$ و یک متغیر هدف $y$ را با یک خط راست مدل می‌کند:
\[
\hat{y} = \beta_0 + \beta_1 x
\]
که در آن $\beta_0$ عرض از مبدأ و $\beta_1$ شیب خط است. این ضرایب معمولاً
با روش \dmterm{\textbf{کمترین مربعات}}{Ordinary Least Squares (OLS)}
تخمین زده می‌شوند؛ یعنی مقادیری برای $\beta_0$ و $\beta_1$ انتخاب
می‌شوند که مجموع مربعات خطای پیش‌بینی (فاصلهٔ عمودی هر نقطهٔ داده تا خط
برازش‌شده) را کمینه کنند.

وقتی بیش از یک ویژگی ورودی داشته باشیم، از
\dmterm{\textbf{رگرسیون خطی چندگانه}}{Multiple Linear Regression}
استفاده می‌شود که رابطه را به‌صورت زیر تعمیم می‌دهد:
\[
\hat{y} = \beta_0 + \beta_1 x_1 + \beta_2 x_2 + \cdots + \beta_n x_n
\]

نکتهٔ مهم در رگرسیون چندگانه، توجه به
\dmterm{\textbf{هم‌خطی چندگانه}}{Multicollinearity} است: اگر دو یا چند
ویژگی ورودی به‌شدت با یکدیگر همبسته باشند، تفسیر ضرایب مدل دشوار و
ناپایدار می‌شود، هرچند ممکن است دقت پیش‌بینی کلی مدل تحت تأثیر جدی قرار
نگیرد.

\section{رگرسیون لجستیک}

با وجود نامش، \dmterm{\textbf{رگرسیون لجستیک}}{Logistic Regression}
عملاً یک الگوریتم \textbf{طبقه‌بندی} است، نه رگرسیون به معنای پیش‌بینی
یک مقدار پیوسته؛ این عنوان از تاریخچهٔ توسعهٔ روش باقی مانده است. رگرسیون
لجستیک برای مسائل طبقه‌بندی دوکلاسه به کار می‌رود و به‌جای پیش‌بینی
مستقیم کلاس، \textbf{احتمال} تعلق به کلاس مثبت را تخمین می‌زند.

ایدهٔ اصلی این است که به‌جای برازش یک خط راست (که می‌تواند مقادیری خارج
از بازهٔ $[0, 1]$ تولید کند و برای احتمال بی‌معنی است)، از
\dmterm{\textbf{تابع سیگموید}}{Sigmoid Function} برای نگاشت خروجی مدل
خطی به بازهٔ $[0, 1]$ استفاده می‌شود:
\[
P(y=1 \mid x) = \frac{1}{1 + e^{-(\beta_0 + \beta_1 x_1 + \cdots + \beta_n x_n)}}
\]

در نهایت، با مقایسهٔ این احتمال با یک آستانه (معمولاً ۰.۵)، تصمیم نهایی
کلاس گرفته می‌شود. به همین دلیل، رگرسیون لجستیک را می‌توان پلی میان
مباحث این فصل و مباحث طبقه‌بندی فصل‌های قبل دانست، و از تمام معیارهای
ارزیابی طبقه‌بندی (ماتریس درهم‌ریختگی، دقت مثبت، فراخوانی، منحنی
\LR{ROC}) که در فصل \ref{ch:classification-advanced} دیدیم، برای آن نیز
استفاده می‌شود.

\section{ارزیابی مدل‌های رگرسیون}

از آنجا که خروجی رگرسیون عددی پیوسته است، معیارهای مبتنی بر ماتریس
درهم‌ریختگی (که برای کلاس‌های گسسته تعریف شده‌اند) در اینجا کاربرد
ندارند. به‌جای آن، معیارهای زیر بر پایهٔ فاصلهٔ میان مقدار پیش‌بینی‌شده
$\hat{y}_i$ و مقدار واقعی $y_i$ تعریف می‌شوند.

\paragraph{میانگین مربعات خطا}
\dmterm{\textbf{میانگین مربعات خطا}}{Mean Squared Error (MSE)} میانگین
مجذور تفاضل پیش‌بینی و مقدار واقعی است:
\[
\text{MSE} = \frac{1}{n} \sum_{i=1}^{n} (y_i - \hat{y}_i)^2
\]
به دلیل به توان دو رساندن خطا، این معیار به خطاهای بزرگ (نقاط پرت) حساسیت
بیشتری نسبت به خطاهای کوچک نشان می‌دهد.

\paragraph{جذر میانگین مربعات خطا}
\dmterm{\textbf{جذر میانگین مربعات خطا}}{Root Mean Squared Error (RMSE)}
با گرفتن جذر از \LR{MSE} به دست می‌آید:
\[
\text{RMSE} = \sqrt{\text{MSE}}
\]
مزیت اصلی \LR{RMSE} نسبت به \LR{MSE} این است که واحد آن با واحد متغیر
هدف یکسان است (مثلاً اگر هدف قیمت به تومان باشد، \LR{RMSE} نیز بر حسب
تومان قابل تفسیر است)، در حالی که واحد \LR{MSE} مجذور واحد اصلی است.

\paragraph{ضریب تعیین}
\dmterm{\textbf{ضریب تعیین}}{Coefficient of Determination ($R^2$)} نشان
می‌دهد چه نسبتی از واریانس متغیر هدف توسط مدل توضیح داده می‌شود:
\[
R^2 = 1 - \frac{\sum_{i=1}^{n} (y_i - \hat{y}_i)^2}{\sum_{i=1}^{n} (y_i - \bar{y})^2}
\]
که در آن $\bar{y}$ میانگین مقادیر واقعی است. مقدار $R^2$ نزدیک به ۱
نشان‌دهندهٔ برازش خوب مدل و مقدار نزدیک به صفر (یا حتی منفی) نشان‌دهندهٔ
عملکرد ضعیف آن است.

\section*{پیاده‌سازی در پایتون}

پیاده‌سازی \texttt{LinearRegression} و \texttt{LogisticRegression} از
\texttt{scikit-learn}، به‌همراه محاسبهٔ \texttt{mean\_squared\_error} و
\texttt{r2\_score}، در بخش «فصل ۵» پیوست الف (\ref{app:python}) آمده
است.

\section*{خلاصهٔ فصل}

در این فصل با رگرسیون به‌عنوان دومین وظیفهٔ اصلی یادگیری با نظارت — این
بار برای پیش‌بینی مقادیر پیوسته — آشنا شدیم. رگرسیون خطی ساده و چندگانه
را با روش کمترین مربعات بررسی کردیم، رگرسیون لجستیک را به‌عنوان پلی
میان رگرسیون و طبقه‌بندی معرفی کردیم، و در پایان با معیارهای ارزیابی
مخصوص مسائل رگرسیون ($\text{MSE}$، $\text{RMSE}$، $R^2$) آشنا شدیم. با
این فصل، بخش یادگیری با نظارت این کتاب به پایان می‌رسد. در فصل بعد، به
دنیای یادگیری بدون نظارت و اولین موضوع آن، یعنی قوانین انجمنی، وارد
می‌شویم.

\section*{تمرین‌ها}

\begin{enumerate}
    \item توضیح دهید چرا رگرسیون لجستیک، با وجود نامش، در دستهٔ
          الگوریتم‌های طبقه‌بندی قرار می‌گیرد نه رگرسیون.
    \item فرض کنید یک مدل رگرسیون خطی برای پیش‌بینی قیمت خانه دارید که
          دو ویژگی ورودی «متراژ» و «تعداد اتاق» به‌شدت با هم همبسته‌اند.
          مشکل احتمالی این وضعیت را توضیح دهید.
    \item چرا معیار \LR{RMSE} برای گزارش نتیجه به یک ذی‌نفع غیرفنی
          معمولاً مناسب‌تر از \LR{MSE} است؟
    \item اگر یک مدل رگرسیون مقدار $R^2 = 0.85$ داشته باشد، این عدد را
          به زبان ساده برای یک مدیر کسب‌وکار توضیح دهید.
    \item تابع سیگموید را رسم کنید (روی کاغذ یا با پایتون) و توضیح دهید
          چرا برای تبدیل خروجی یک مدل خطی به احتمال مناسب است.
\end{enumerate}


% ---------- بخش سوم: یادگیری بدون نظارت ----------
\part{یادگیری بدون نظارت}
\chapter{قوانین انجمنی}
\label{ch:association}

\section*{اهداف این فصل}
\begin{itemize}
    \item درک مفهوم قوانین انجمنی و کاربرد آن در تحلیل سبد خرید
    \item تسلط بر معیارهای Support، Confidence و Lift
    \item یادگیری الگوریتم Apriori برای یافتن مجموعه‌های پرتکرار
    \item آشنایی با FP-Growth به‌عنوان روشی کارآمدتر
\end{itemize}

\section{مقدمه: از یادگیری با نظارت به بدون نظارت}

در چهار فصل گذشته، با وظایف \textbf{یادگیری با نظارت} (طبقه‌بندی و
رگرسیون) آشنا شدیم که در آن‌ها داده برچسب‌دار بود. از این فصل به بعد،
وارد دنیای \dmterm{\textbf{یادگیری بدون نظارت}}{Unsupervised Learning}
می‌شویم، جایی که هیچ برچسب از پیش‌تعیین‌شده‌ای وجود ندارد و هدف، کشف
ساختار یا الگوهای پنهان در خود داده است.

\dmterm{\textbf{قوانین انجمنی}}{Association Rules} یکی از قدیمی‌ترین و
شناخته‌شده‌ترین وظایف یادگیری بدون نظارت است. مثال کلاسیک آن
\dmterm{\textbf{تحلیل سبد خرید}}{Market Basket Analysis} است: با بررسی
میلیون‌ها رسید خرید یک فروشگاه، آیا می‌توان الگوهایی مانند «مشتریانی که
نان و کره می‌خرند، با احتمال بالا شیر هم می‌خرند» را به‌صورت خودکار کشف
کرد؟ چنین دانشی می‌تواند در چیدمان قفسه‌ها، طراحی تخفیف‌های ترکیبی و
سیستم‌های پیشنهاددهنده استفاده شود.

\section{مفاهیم پایه}

فرض کنید $I = \{i_1, i_2, \ldots, i_m\}$ مجموعهٔ همهٔ کالاهای موجود در
یک فروشگاه باشد و هر \dmterm{تراکنش}{Transaction} $T$ زیرمجموعه‌ای از
$I$ باشد (یعنی سبد خرید یک مشتری). یک \dmterm{\textbf{قانون انجمنی}}{Association
Rule} به‌صورت $A \Rightarrow B$ نوشته می‌شود، که در آن $A$ و $B$ دو
مجموعهٔ کالای مجزا از هم هستند (مثلاً $A = \{$نان، کره$\}$ و
$B = \{$شیر$\}$).

سه معیار زیر برای سنجش کیفیت یک قانون انجمنی به کار می‌روند:

\paragraph{پشتیبانی}
\dmterm{\textbf{پشتیبانی}}{Support} نشان می‌دهد چه نسبتی از کل تراکنش‌ها
شامل هر دو مجموعهٔ $A$ و $B$ با هم هستند:
\[
\text{Support}(A \Rightarrow B) = \frac{|\{T : A \cup B \subseteq T\}|}{|\text{کل تراکنش‌ها}|}
\]
پشتیبانی نشان‌دهندهٔ \textbf{شیوع} یک الگو در کل داده است؛ قاعده‌ای با
پشتیبانی بسیار کم، حتی اگر جالب باشد، ممکن است از نظر تجاری اهمیت کمی
داشته باشد چون به‌ندرت رخ می‌دهد.

\paragraph{اطمینان}
\dmterm{\textbf{اطمینان}}{Confidence} نشان می‌دهد از میان تراکنش‌هایی که
$A$ در آن‌ها وجود دارد، چه نسبتی $B$ را نیز شامل می‌شوند:
\[
\text{Confidence}(A \Rightarrow B) = \frac{\text{Support}(A \cup B)}{\text{Support}(A)}
\]
اطمینان معادل احتمال شرطی $P(B \mid A)$ است.

\paragraph{لیفت}
مشکل اطمینان این است که اگر $B$ در بین همهٔ تراکنش‌ها بسیار رایج باشد
(مثلاً اکثر مشتریان به هر حال نان می‌خرند)، اطمینان بالا می‌تواند
گمراه‌کننده باشد و نشان‌دهندهٔ یک رابطهٔ واقعی نباشد. برای رفع این مشکل،
از معیار \dmterm{\textbf{لیفت}}{Lift} استفاده می‌شود:
\[
\text{Lift}(A \Rightarrow B) = \frac{\text{Confidence}(A \Rightarrow B)}{\text{Support}(B)}
\]
اگر $\text{Lift} > 1$، حضور $A$ احتمال حضور $B$ را افزایش می‌دهد (رابطهٔ
مثبت). اگر $\text{Lift} \approx 1$، دو رویداد تقریباً مستقل‌اند (اطمینان
بالا صرفاً به دلیل رواج بالای $B$ بوده، نه یک رابطهٔ واقعی). اگر
$\text{Lift} < 1$، حضور $A$ احتمال حضور $B$ را کاهش می‌دهد (رابطهٔ
منفی).

\section{الگوریتم Apriori}

چالش اصلی در یافتن قوانین انجمنی، تعداد نمایی زیرمجموعه‌های ممکن از
کالاهاست: برای $n$ کالا، $2^n$ زیرمجموعهٔ ممکن وجود دارد که بررسی همهٔ
آن‌ها برای مجموعه داده‌های واقعی (با هزاران کالا) از نظر محاسباتی
غیرممکن است. \dmterm{\textbf{الگوریتم} \LR{Apriori}}{Apriori Algorithm}
با استفاده از یک ویژگی هوشمندانه به نام
\dmterm{\textbf{اصل Apriori}}{Apriori Principle} این فضای جستجو را
به‌شدت کاهش می‌دهد:

\begin{quote}
اگر یک مجموعه کالا پرتکرار\footnote{\LR{Frequent}} نباشد، هیچ‌یک از
ابرمجموعه‌های آن نیز نمی‌توانند پرتکرار باشند.
\end{quote}

به بیان دیگر، اگر مجموعهٔ $\{$نان، کره$\}$ پشتیبانی کافی نداشته باشد،
نیازی به بررسی $\{$نان، کره، شیر$\}$ نیست، چون پشتیبانی آن هرگز نمی‌تواند
از پشتیبانی زیرمجموعه‌اش بیشتر باشد.

\subsection{تولید مجموعه‌های پرتکرار}

الگوریتم Apriori به‌صورت تکرارشونده و لایه‌به‌لایه عمل می‌کند:

\begin{enumerate}
    \item ابتدا پشتیبانی همهٔ مجموعه‌های تک‌عضوی محاسبه می‌شود و آن‌هایی
          که از یک آستانهٔ حداقلی\footnote{\LR{Minimum Support}} پایین‌ترند،
          حذف می‌شوند.
    \item از ترکیب مجموعه‌های پرتکرار باقی‌مانده، مجموعه‌های دوعضوی
          کاندیدا ساخته می‌شوند و دوباره پشتیبانی آن‌ها بررسی می‌شود.
    \item این فرآیند برای مجموعه‌های سه‌عضوی، چهارعضوی و به همین ترتیب
          ادامه می‌یابد تا مجموعهٔ کاندیدای جدیدی باقی نماند.
\end{enumerate}

\subsection{استخراج قوانین}

پس از یافتن همهٔ مجموعه‌های پرتکرار، برای هر مجموعه، تمام تفکیک‌های
ممکن آن به دو بخش $A$ و $B$ بررسی می‌شود و اطمینان هر قانون محاسبه
می‌گردد. قوانینی که اطمینان آن‌ها از یک آستانهٔ حداقلی
\footnote{\LR{Minimum Confidence}} بیشتر باشد، به‌عنوان قوانین نهایی
گزارش می‌شوند.

\paragraph{ضعف اصلی Apriori}
ضعف عملی این الگوریتم، نیاز به بارها اسکن کردن کل مجموعه داده (یک بار
به ازای هر سطح از اندازهٔ مجموعه کاندیدا) است که برای داده‌های بسیار
بزرگ، هزینهٔ زمانی و حافظه‌ای بالایی دارد.

\section{الگوریتم FP-Growth}

\dmterm{\textbf{الگوریتم} \LR{FP-Growth}}{Frequent Pattern Growth} برای
رفع ضعف اصلی Apriori طراحی شده است. ایدهٔ اصلی آن، فشرده‌سازی کل
مجموعه داده در یک ساختار درختی به نام
\dmterm{\textbf{درخت} \LR{FP}}{FP-Tree} است، به‌گونه‌ای که تنها با دو
بار اسکن کل داده (به‌جای اسکن‌های مکرر Apriori)، می‌توان همهٔ مجموعه‌های
پرتکرار را با پیمایش این درخت استخراج کرد. به همین دلیل، FP-Growth در
عمل معمولاً به‌مراتب سریع‌تر از Apriori است، هرچند پیاده‌سازی آن پیچیده‌تر
است. برای اهداف آموزشی این کتاب، درک اصل Apriori اولویت دارد؛ FP-Growth
را می‌توان نسخهٔ بهینه‌شدهٔ همان ایده در نظر گرفت.

\section{ارزیابی و انتخاب قوانین جالب}

الگوریتم‌های یافتن قوانین انجمنی معمولاً تعداد بسیار زیادی قانون تولید
می‌کنند که اغلب آن‌ها بی‌فایده یا بدیهی‌اند (مانند «مشتریانی که نان
می‌خرند، احتمالاً نان می‌خواهند»). برای انتخاب قوانین واقعاً
\dmterm{\textbf{جالب}}{Interesting}، معیار لیفت (که پیش‌تر دیدیم) نقش
کلیدی دارد؛ قوانینی با لیفت نزدیک به ۱ معمولاً کنار گذاشته می‌شوند، حتی
اگر پشتیبانی و اطمینان بالایی داشته باشند. علاوه بر این، قضاوت متخصص
حوزه (مثلاً یک کارشناس بازاریابی که می‌داند کدام رابطه از قبل شناخته‌شده
و کدام واقعاً جدید است) نقش تکمیلی مهمی در فیلتر نهایی قوانین ایفا
می‌کند.

\section*{پیاده‌سازی در پایتون}

پیاده‌سازی الگوریتم Apriori و استخراج قوانین با کتابخانهٔ
\texttt{mlxtend} (توابع \texttt{apriori} و \texttt{association\_rules})
در بخش «فصل ۶» پیوست الف (\ref{app:python}) آمده است.

\section*{خلاصهٔ فصل}

در این فصل، اولین وظیفهٔ یادگیری بدون نظارت این کتاب، یعنی قوانین
انجمنی، را بررسی کردیم. با سه معیار کلیدی — پشتیبانی، اطمینان و لیفت —
آشنا شدیم و دیدیم چگونه اصل Apriori فضای جستجوی نمایی مجموعه‌های کالا
را قابل مدیریت می‌کند. همچنین به‌طور مختصر FP-Growth را به‌عنوان
جایگزینی کارآمدتر معرفی کردیم. در فصل بعد، به دومین وظیفهٔ اصلی یادگیری
بدون نظارت، یعنی خوشه‌بندی، می‌پردازیم.

\section*{تمرین‌ها}

\begin{enumerate}
    \item از میان ۱۰۰۰ تراکنش یک فروشگاه، ۲۰۰ تراکنش شامل «نان» و ۵۰
          تراکنش شامل هم «نان» و هم «کره» هستند. پشتیبانی و اطمینان قانون
          $\{$نان$\} \Rightarrow \{$کره$\}$ را محاسبه کنید.
    \item اگر پشتیبانی $\{$کره$\}$ در کل داده ۱۵ درصد باشد، لیفت قانون
          تمرین قبل را محاسبه کنید و نتیجه را تفسیر کنید (آیا نان و کره
          رابطهٔ مثبت، منفی یا مستقل دارند؟).
    \item اصل Apriori را با کلمات خودتان توضیح دهید و بگویید چرا این اصل
          فضای جستجو را به‌شدت کاهش می‌دهد.
    \item چرا یک قانون می‌تواند اطمینان بسیار بالایی داشته باشد (مثلاً ۹۵
          درصد) اما از نظر تجاری بی‌ارزش باشد؟ با یک مثال توضیح دهید.
    \item مزیت اصلی FP-Growth نسبت به Apriori از نظر تعداد بار اسکن کردن
          داده چیست؟
\end{enumerate}

\chapter{خوشه‌بندی}
\label{ch:clustering}

\section*{اهداف این فصل}
\begin{itemize}
    \item درک مفهوم خوشه‌بندی و تفاوت آن با طبقه‌بندی
    \item یادگیری الگوریتم k-means و نحوهٔ انتخاب مقدار $k$
    \item آشنایی با خوشه‌بندی سلسله‌مراتبی و روش‌های پیوند
    \item آشنایی با DBSCAN برای خوشه‌های با شکل دلخواه
    \item آشنایی با معیارهای ارزیابی خوشه‌بندی
\end{itemize}

\section{مقدمه‌ای بر خوشه‌بندی}

\dmterm{\textbf{خوشه‌بندی}}{Clustering} دومین وظیفهٔ اصلی یادگیری بدون
نظارت است. هدف آن، گروه‌بندی نمونه‌ها به‌گونه‌ای است که نمونه‌های درون
یک گروه (\dmterm{خوشه}{Cluster}) شباهت زیادی به هم داشته باشند و
نمونه‌های دو گروه مختلف، شباهت کمی به یکدیگر داشته باشند — بدون آنکه
هیچ برچسب از پیش‌تعیین‌شده‌ای در اختیار باشد. این وظیفه با طبقه‌بندی
تفاوت اساسی دارد: در طبقه‌بندی، کلاس‌ها از پیش مشخص‌اند و مدل یاد
می‌گیرد نمونهٔ جدید به کدام کلاس تعلق دارد؛ در خوشه‌بندی، حتی تعداد و
ماهیت گروه‌ها نیز باید از دل خود داده کشف شود.

کاربردهای رایج خوشه‌بندی عبارت‌اند از بخش‌بندی مشتریان بر اساس رفتار
خرید، فشرده‌سازی تصویر، و به‌عنوان گامی مقدماتی پیش از تحلیل‌های دیگر
(مثلاً یافتن زیرگروه‌های همگن در یک جمعیت پیش از مدل‌سازی جداگانه برای
هر زیرگروه).

\section{الگوریتم k-means}

\dmterm{\textbf{الگوریتم} \LR{k-means}}{k-means Algorithm} پرکاربردترین
الگوریتم خوشه‌بندی است. ایدهٔ آن ساده است: داده به $k$ خوشه تقسیم
می‌شود، به‌گونه‌ای که هر نمونه به نزدیک‌ترین
\dmterm{مرکز خوشه}{Centroid} تعلق داشته باشد. الگوریتم به‌صورت
تکرارشونده عمل می‌کند:

\begin{enumerate}
    \item ابتدا $k$ مرکز خوشه به‌طور تصادفی انتخاب می‌شوند.
    \item هر نمونه به نزدیک‌ترین مرکز خوشه (بر اساس فاصلهٔ اقلیدسی)
          نسبت داده می‌شود.
    \item مرکز هر خوشه، به‌عنوان میانگین همهٔ نمونه‌های آن خوشه، دوباره
          محاسبه می‌شود.
    \item مراحل ۲ و ۳ تا زمانی تکرار می‌شوند که مراکز خوشه دیگر تغییر
          معناداری نکنند (همگرایی).
\end{enumerate}

هدف k-means کمینه‌کردن مجموع مربعات فاصلهٔ درون‌خوشه‌ای است
\footnote{\LR{Within-Cluster Sum of Squares (WCSS)}}. نکات مهم عملی:

\begin{itemize}
    \item چون انتخاب اولیهٔ مراکز تصادفی است، نتیجهٔ نهایی می‌تواند به
          این انتخاب حساس باشد؛ در عمل الگوریتم چند بار با مقداردهی
          اولیهٔ متفاوت اجرا و بهترین نتیجه انتخاب می‌شود.
    \item k-means برای خوشه‌های با شکل تقریباً کروی و اندازهٔ مشابه بهتر
          عمل می‌کند و برای خوشه‌های با شکل دلخواه یا اندازهٔ بسیار
          متفاوت مناسب نیست (بخش \ref{sec:dbscan} را ببینید).
    \item از آنجا که بر پایهٔ فاصله عمل می‌کند، نرمال‌سازی ویژگی‌ها
          (فصل \ref{ch:data}) پیش از اجرا ضروری است.
\end{itemize}

\subsection{انتخاب مقدار k}

برخلاف طبقه‌بندی که تعداد کلاس‌ها از قبل مشخص است، در k-means باید
مقدار $k$ را از قبل تعیین کنیم. دو روش رایج برای این کار عبارت‌اند از:

\paragraph{روش آرنج}
در \dmterm{\textbf{روش آرنج}}{Elbow Method}، الگوریتم برای مقادیر
مختلف $k$ اجرا می‌شود و مقدار \LR{WCSS} برای هر $k$ رسم می‌شود. با
افزایش $k$، این مقدار همواره کاهش می‌یابد، اما نرخ کاهش معمولاً در یک
نقطهٔ خاص به‌شدت کم می‌شود و نموداری شبیه آرنج بازو ایجاد می‌کند؛ آن
نقطه، انتخاب مناسبی برای $k$ است.

\paragraph{ضریب سیلوئت}
\dmterm{\textbf{ضریب سیلوئت}}{Silhouette Coefficient} برای هر نمونه،
میزان شباهت آن به خوشهٔ خودش را در مقایسه با نزدیک‌ترین خوشهٔ دیگر
می‌سنجد و مقداری بین $-1$ تا $1$ می‌دهد؛ مقدار نزدیک به $1$ نشان‌دهندهٔ
خوشه‌بندی مناسب است. میانگین این ضریب برای همهٔ نمونه‌ها، معیاری کلی
برای مقایسهٔ کیفیت خوشه‌بندی با مقادیر مختلف $k$ به دست می‌دهد و
برخلاف روش آرنج، نیازی به قضاوت چشمی ندارد.

\section{خوشه‌بندی سلسله‌مراتبی}

برخلاف k-means که تعداد خوشه‌ها را از پیش می‌خواهد،
\dmterm{\textbf{خوشه‌بندی سلسله‌مراتبی}}{Hierarchical Clustering}
ساختاری درختی از خوشه‌ها در سطوح مختلف می‌سازد که با نموداری به نام
\dmterm{\textbf{دندروگرام}}{Dendrogram} نمایش داده می‌شود؛ تعداد نهایی
خوشه‌ها را می‌توان بعداً با «بریدن» این درخت در ارتفاع دلخواه تعیین
کرد. رایج‌ترین رویکرد،
\dmterm{\textbf{رویکرد تجمیعی}}{Agglomerative Approach} است که به‌صورت
پایین به بالا عمل می‌کند: در ابتدا هر نمونه یک خوشهٔ جداگانه است، و در
هر مرحله، دو خوشهٔ نزدیک‌تر به هم ادغام می‌شوند تا در نهایت یک خوشهٔ
واحد باقی بماند.

\subsection{روش‌های پیوند}

معیار «نزدیکی» میان دو خوشه (نه دو نمونهٔ منفرد) با
\dmterm{\textbf{روش‌های پیوند}}{Linkage Methods} تعریف می‌شود:

\begin{itemize}
    \item \dmterm{\textbf{پیوند تکی}}{Single Linkage}: کمترین فاصله میان
          هر عضو یک خوشه با هر عضو خوشهٔ دیگر.
    \item \dmterm{\textbf{پیوند کامل}}{Complete Linkage}: بیشترین فاصله
          میان اعضای دو خوشه.
    \item \dmterm{\textbf{پیوند میانگین}}{Average Linkage}: میانگین
          فاصلهٔ همهٔ جفت‌های ممکن میان اعضای دو خوشه.
\end{itemize}

پیوند تکی می‌تواند به پدیدهٔ \dmterm{زنجیره‌ای‌شدن}{Chaining} (تشکیل
خوشه‌های کشیده و نامناسب) منجر شود، در حالی که پیوند کامل معمولاً خوشه‌های
فشرده‌تر و کروی‌تری تولید می‌کند.

\section{DBSCAN}
\label{sec:dbscan}

\dmterm{\textbf{الگوریتم} \LR{DBSCAN}}{Density-Based Spatial Clustering
of Applications with Noise} برخلاف دو روش قبل، خوشه‌ها را بر اساس
\textbf{چگالی} نقاط تعریف می‌کند، نه فاصله تا یک مرکز یا ساختار
سلسله‌مراتبی. این ویژگی به آن اجازه می‌دهد خوشه‌هایی با شکل دلخواه
(نه‌فقط کروی) را شناسایی کند و به‌طور طبیعی، نقاطی که در هیچ ناحیهٔ
پرچگالی قرار نمی‌گیرند را به‌عنوان \textbf{نویز} یا داده پرت علامت‌گذاری
کند (که ارتباط مستقیمی با فصل \ref{ch:outliers} دارد).

DBSCAN دو پارامتر اصلی دارد: شعاع همسایگی ($\varepsilon$) و حداقل تعداد
نقاط لازم در آن همسایگی برای تشکیل یک ناحیهٔ پرچگالی\footnote{\LR{MinPts}}.
نقاطی که این حداقل تعداد همسایه را در شعاع خود
داشته باشند \dmterm{نقاط هسته‌ای}{Core Points} نامیده می‌شوند و خوشه‌ها
با اتصال نقاط هسته‌ای مجاور به یکدیگر شکل می‌گیرند.

\section{ارزیابی خوشه‌بندی}

از آنجا که در خوشه‌بندی برچسب واقعی وجود ندارد، ارزیابی آن دشوارتر از
طبقه‌بندی است. رایج‌ترین رویکرد، استفاده از همان
\dmterm{ضریب سیلوئت}{Silhouette Coefficient} است که در بخش انتخاب $k$
دیدیم؛ این معیار می‌تواند برای مقایسهٔ کیفیت هر خوشه‌بندی (فارغ از
الگوریتم تولیدکننده) به کار رود. در مواردی که برچسب واقعی (برای اهداف
پژوهشی) در دسترس باشد، معیارهایی مانند
\dmterm{\textbf{شاخص رند تعدیل‌شده}}{Adjusted Rand Index} نیز برای
مقایسهٔ خوشه‌بندی به‌دست‌آمده با برچسب واقعی استفاده می‌شوند.

\section*{پیاده‌سازی در پایتون}

پیاده‌سازی \texttt{KMeans}، \texttt{AgglomerativeClustering} و
\texttt{DBSCAN} از \texttt{scikit-learn}، به‌همراه محاسبهٔ ضریب سیلوئت
و رسم دندروگرام، در بخش «فصل ۷» پیوست الف (\ref{app:python}) آمده است.

\section*{خلاصهٔ فصل}

در این فصل با خوشه‌بندی به‌عنوان دومین وظیفهٔ اصلی یادگیری بدون نظارت
آشنا شدیم. سه رویکرد اصلی را بررسی کردیم: k-means (مبتنی بر فاصله تا
مرکز)، خوشه‌بندی سلسله‌مراتبی (مبتنی بر ساختار درختی و روش‌های پیوند)
و DBSCAN (مبتنی بر چگالی). همچنین با روش آرنج و ضریب سیلوئت برای انتخاب
تعداد خوشه‌ها و ارزیابی کیفیت خوشه‌بندی آشنا شدیم. در فصل بعد، به
موضوع نزدیک به خوشه‌بندی، یعنی تشخیص داده‌های پرت، می‌پردازیم.

\section*{تمرین‌ها}

\begin{enumerate}
    \item تفاوت اصلی خوشه‌بندی با طبقه‌بندی را از نظر نوع داده مورد نیاز
          (برچسب‌دار در برابر بدون برچسب) توضیح دهید.
    \item چرا نتیجهٔ k-means می‌تواند در دو اجرای مختلف با یک مجموعه
          داده ثابت، متفاوت باشد؟ چگونه می‌توان این مشکل را کاهش داد؟
    \item تفاوت پیوند تکی و پیوند کامل را با یک مثال از شکل خوشه‌های
          حاصل توضیح دهید.
    \item چرا DBSCAN می‌تواند خوشه‌هایی با شکل دلخواه را شناسایی کند در
          حالی که k-means معمولاً خوشه‌های کروی تولید می‌کند؟
    \item توضیح دهید ضریب سیلوئت چگونه هم برای انتخاب مقدار $k$ در
          k-means و هم برای مقایسهٔ کلی کیفیت خوشه‌بندی به کار می‌رود.
\end{enumerate}

\chapter{تشخیص داده‌های پرت}
\label{ch:outliers}

\section*{اهداف این فصل}
\begin{itemize}
    \item درک مفهوم داده پرت و تفاوت آن با نویز
    \item آشنایی با روش‌های آماری ساده برای تشخیص داده پرت
    \item آشنایی با روش‌های مبتنی بر فاصله و چگالی
    \item آشنایی با روش‌های مبتنی بر یادگیری ماشین
\end{itemize}

\section{تعریف و انواع داده‌های پرت}

همان‌طور که در فصل \ref{ch:data} اشاره شد،
\dmterm{\textbf{داده پرت}}{Outlier} نمونه‌ای است که به‌طور قابل‌توجهی با
بقیهٔ داده تفاوت دارد. برخلاف نویز که معمولاً خطای بی‌معنی و قابل‌حذف
است، داده پرت می‌تواند خودِ هدف اصلی تحلیل باشد: تراکنش متقلبانه در
داده‌های بانکی، رفتار غیرعادی یک حسگر پیش از خرابی تجهیزات صنعتی، یا
یک بیمار با علائم نادر و نگران‌کننده. به همین دلیل، این وظیفه گاهی
\dmterm{\textbf{تشخیص ناهنجاری}}{Anomaly Detection} نیز نامیده می‌شود.

داده‌های پرت را می‌توان به سه دسته تقسیم کرد:

\begin{itemize}
    \item \dmterm{\textbf{پرت سراسری}}{Global Outlier}: نمونه‌ای که در
          مقایسه با \textbf{کل} مجموعه داده غیرعادی است (مثلاً دمای بدن
          ۴۲ درجه).
    \item \dmterm{\textbf{پرت زمینه‌ای}}{Contextual Outlier}: نمونه‌ای
          که تنها در یک زمینهٔ خاص غیرعادی است (مثلاً دمای هوای ۳۵ درجه
          در زمستان غیرعادی است، اما در تابستان طبیعی است).
    \item \dmterm{\textbf{پرت جمعی}}{Collective Outlier}: مجموعه‌ای از
          نمونه‌ها که هرکدام به‌تنهایی عادی به نظر می‌رسند، اما رفتار
          جمعی آن‌ها غیرعادی است (مثلاً یک الگوی ترافیک شبکه که نشانهٔ
          حملهٔ سایبری است).
\end{itemize}

\section{روش‌های آماری}

ساده‌ترین رویکرد برای تشخیص داده پرت، استفاده از فرضیات آماری دربارهٔ
توزیع داده است. برای داده‌ای که تقریباً توزیع نرمال دارد، رایج‌ترین
روش استفاده از \dmterm{\textbf{امتیاز} \LR{Z}}{Z-Score} است (که پیش‌تر
در فصل \ref{ch:data} برای استانداردسازی دیدیم): نمونه‌ای که قدرمطلق
امتیاز $Z$ آن بیش از یک آستانه (معمولاً ۲.۵ یا ۳) باشد، به‌عنوان داده
پرت علامت‌گذاری می‌شود.

روش دیگر که وابستگی کمتری به فرض نرمال بودن توزیع دارد، استفاده از
\dmterm{\textbf{دامنهٔ میان‌چارکی}}{Interquartile Range (IQR)} است: هر
مقداری که خارج از بازهٔ
$[Q_1 - 1.5 \times IQR,\ Q_3 + 1.5 \times IQR]$ قرار گیرد (که در آن
$Q_1$ و $Q_3$ به‌ترتیب چارک اول و سوم داده‌اند)، به‌عنوان پرت در نظر
گرفته می‌شود. این روش پایهٔ نمودار جعبه‌ای است\footnote{\LR{Box Plot}}.

محدودیت اصلی روش‌های آماری این است که معمولاً فقط برای یک ویژگی در یک
زمان مناسب‌اند (تک‌متغیره) و به‌سختی به داده‌های چندبعدی با روابط
پیچیده میان ویژگی‌ها تعمیم می‌یابند.

\section{روش‌های مبتنی بر فاصله و چگالی}

روش‌های مبتنی بر فاصله، داده پرت را نمونه‌ای تعریف می‌کنند که همسایگان
نزدیک کمی دارد یا فاصلهٔ آن تا نزدیک‌ترین همسایه‌هایش زیاد است — رویکردی
مشابه چیزی که در DBSCAN (فصل \ref{ch:clustering}) دیدیم، که در آن نقاط
کم‌چگالی به‌طور طبیعی به‌عنوان نویز شناسایی می‌شدند.

روش دقیق‌تری که چگالی \textbf{نسبی} را در نظر می‌گیرد،
\dmterm{\textbf{عامل پرت محلی}}{Local Outlier Factor (LOF)} نام دارد.
مزیت اصلی LOF نسبت به روش‌های ساده‌تر این است که چگالی هر نقطه را نسبت
به چگالی همسایگانش می‌سنجد، نه نسبت به یک آستانهٔ سراسری ثابت؛ این
ویژگی به آن اجازه می‌دهد داده‌های پرت را حتی در مجموعه‌هایی که چگالی
خوشه‌های مختلف با هم تفاوت دارد (که در آن‌ها روش‌های سراسری‌تر دچار خطا
می‌شوند)، به‌درستی شناسایی کند.

\section{روش‌های مبتنی بر یادگیری}

\dmterm{\textbf{جنگل ایزوله}}{Isolation Forest} یکی از مؤثرترین
روش‌های مدرن تشخیص داده پرت است. ایدهٔ اصلی آن برخلاف روش‌های قبلی که
بر پایهٔ فاصله یا چگالی «عادی بودن» را می‌سنجیدند، مستقیماً بر پایهٔ
سهولت \textbf{ایزوله‌کردن} یک نمونه است: با ساخت درخت‌های تصادفی که
به‌طور مکرر داده را بر اساس ویژگی‌ها و مقادیر تصادفی تقسیم می‌کنند،
نمونه‌های پرت — که از بقیهٔ داده متمایزند — معمولاً با تعداد تقسیم بسیار
کمتری ایزوله می‌شوند (یعنی در عمق کمتری از درخت به تنهایی قرار
می‌گیرند)، در حالی که نمونه‌های عادی برای ایزوله شدن به تقسیم‌های بیشتری
نیاز دارند. میانگین عمق ایزوله‌سازی روی چند درخت تصادفی، امتیاز پرت
بودن هر نمونه را می‌دهد.

مزیت اصلی جنگل ایزوله نسبت به روش‌های مبتنی بر فاصله (مانند LOF)، کارایی
محاسباتی بهتر آن روی داده‌های با ابعاد بالا و حجم زیاد است، زیرا نیازی
به محاسبهٔ فاصلهٔ جفت‌به‌جفت میان همهٔ نمونه‌ها ندارد.

\section*{پیاده‌سازی در پایتون}

پیاده‌سازی \texttt{IsolationForest} و \texttt{LocalOutlierFactor} از
\texttt{scikit-learn}، به‌همراه محاسبهٔ امتیاز $Z$ و دامنهٔ میان‌چارکی
با \texttt{numpy}/\texttt{pandas}، در بخش «فصل ۸» پیوست الف
(\ref{app:python}) آمده است.

\section*{خلاصهٔ فصل}

در این فصل با مفهوم داده پرت، سه نوع اصلی آن (سراسری، زمینه‌ای، جمعی)
و طیفی از روش‌های تشخیص آن — از روش‌های آماری ساده (امتیاز $Z$، دامنهٔ
میان‌چارکی) تا روش‌های مبتنی بر چگالی (\LR{LOF}) و روش‌های مدرن مبتنی
بر یادگیری (جنگل ایزوله) — آشنا شدیم. با این فصل، بخش «یادگیری بدون
نظارت» این کتاب تکمیل می‌شود. در فصل بعد، به‌طور مختصر به دو حوزهٔ
تکمیلی، یعنی داده‌کاوی متن و وب، می‌پردازیم.

\section*{تمرین‌ها}

\begin{enumerate}
    \item برای هر یک از مثال‌های زیر مشخص کنید که به کدام نوع داده پرت
          (سراسری، زمینه‌ای یا جمعی) نزدیک‌تر است: (الف) خرید ناگهانی
          و بسیار بزرگ روی یک کارت اعتباری در نیمهٔ شب، (ب) یک روز با
          دمای ۳۰ درجه در وسط زمستان، (ج) یک سری کوتاه از درخواست‌های
          شبکه با الگوی غیرعادی که نشانهٔ حمله است.
    \item چرا روش‌های آماری تک‌متغیره (مانند امتیاز $Z$) برای داده‌های
          چندبعدی با روابط پیچیده میان ویژگی‌ها کافی نیستند؟
    \item تفاوت اصلی \LR{LOF} با روش‌های ساده‌تر مبتنی بر فاصله (مانند
          فاصله تا $k$اُمین همسایهٔ نزدیک) را در مواجهه با خوشه‌هایی با
          چگالی متفاوت توضیح دهید.
    \item ایدهٔ اصلی جنگل ایزوله را در یک یا دو جمله توضیح دهید و بگویید
          چرا این ایده برای داده‌های با ابعاد بالا کارآمد است.
    \item یک سناریوی واقعی (غیر از مثال‌های این فصل) پیدا کنید که در آن
          تشخیص داده پرت اهمیت عملی داشته باشد و توضیح دهید کدام‌یک از
          روش‌های این فصل برای آن مناسب‌تر است.
\end{enumerate}


% ---------- بخش چهارم: موضوعات تکمیلی ----------
\part{موضوعات تکمیلی}
\chapter{مقدمه‌ای بر داده‌کاوی متن و وب}
\label{ch:text-web}

\section*{اهداف این فصل}
\begin{itemize}
    \item آشنایی با گام‌های اصلی پیش‌پردازش متن
    \item یادگیری بازنمایی متن به‌صورت عددی (Bag-of-Words و TF-IDF)
    \item آشنایی مقدماتی با کاربرد طبقه‌بندی و خوشه‌بندی روی متن
    \item آشنایی کلی با مفاهیم پایهٔ داده‌کاوی وب
\end{itemize}

\section{چرا متن به رفتار متفاوتی نیاز دارد؟}

تمام الگوریتم‌هایی که تاکنون دیدیم (درخت تصمیم، k-means، k-NN و
مانند آن‌ها) به یک بازنمایی عددی و با ابعاد ثابت از هر نمونه نیاز
دارند. اما متن خام — یک ایمیل، یک نظر مشتری، یک مقالهٔ خبری — رشته‌ای
از کلمات با طول متغیر است که مستقیماً قابل استفاده در این الگوریتم‌ها
نیست. بنابراین اولین گام در هر پروژهٔ
\dmterm{\textbf{داده‌کاوی متن}}{Text Mining}، تبدیل متن به یک بازنمایی
عددی مناسب است.

\section{پیش‌پردازش متن}

پیش از تبدیل متن به بردار عددی، معمولاً چند گام پیش‌پردازش زیر انجام
می‌شود:

\begin{itemize}
    \item \dmterm{\textbf{نشانه‌گذاری}}{Tokenization}: شکستن متن به
          واحدهای کوچک‌تر، معمولاً کلمات یا زیرکلمات\footnote{\LR{Tokens}}.
    \item \dmterm{\textbf{حذف ایست‌واژه‌ها}}{Stop-word Removal}: حذف
          کلمات بسیار پرتکرار و کم‌اطلاعات مانند حروف اضافه و ضمایر
          (مثلاً «را»، «به»، «از» در فارسی، یا «the»، «is» در انگلیسی)
          که معمولاً برای تشخیص موضوع متن مفید نیستند.
    \item \dmterm{\textbf{ریشه‌یابی}}{Stemming} یا
          \dmterm{\textbf{بن‌واژه‌یابی}}{Lemmatization}: تبدیل کلمات
          مختلف با ریشهٔ مشترک به یک شکل واحد (مثلاً «رفتن»، «می‌رود»،
          «رفت» به یک ریشهٔ مشترک)، تا الگوریتم آن‌ها را یک مفهوم واحد
          در نظر بگیرد.
    \item نرمال‌سازی متن: تبدیل به حروف کوچک (در زبان‌های لاتین)، حذف
          علائم نگارشی و اعداد نامرتبط.
\end{itemize}

\section{بازنمایی متن به‌صورت عددی}

\subsection{کیسهٔ کلمات}

ساده‌ترین روش بازنمایی متن،
\dmterm{\textbf{کیسهٔ کلمات}}{Bag-of-Words (BoW)} است: هر سند با یک
بردار به طول واژگان کل مجموعه داده نمایش داده می‌شود که هر درایهٔ آن،
تعداد تکرار یک کلمهٔ خاص در آن سند است. نام «کیسه» به این دلیل است که
این روش کاملاً \textbf{ترتیب} کلمات را نادیده می‌گیرد؛ برای این روش،
جملهٔ «سگ گربه را دنبال کرد» و «گربه سگ را دنبال کرد» کاملاً یکسان
دیده می‌شوند.

\subsection{TF-IDF}

مشکل اصلی کیسهٔ کلمات این است که به کلمات بسیار پرتکرار (که لزوماً
اطلاعات زیادی دربارهٔ موضوع سند نمی‌دهند) وزن بیش از حد می‌دهد.
\dmterm{\textbf{\LR{TF-IDF}}}{Term Frequency-Inverse Document Frequency}
این مشکل را با ترکیب دو عامل حل می‌کند:

\begin{itemize}
    \item \dmterm{\textbf{فراوانی واژه}}{Term Frequency (TF)}: تعداد
          تکرار یک کلمه در یک سند خاص (هرچه بیشتر، وزن بیشتر).
    \item \dmterm{\textbf{فراوانی معکوس سند}}{Inverse Document Frequency
          (IDF)}: هرچه یک کلمه در تعداد \textbf{کمتری} از اسناد کل
          مجموعه ظاهر شود، وزن بیشتری می‌گیرد؛ کلماتی که تقریباً در همهٔ
          اسناد تکرار می‌شوند (مانند حروف اضافهٔ باقی‌مانده)، وزن بسیار
          کمی می‌گیرند.
\end{itemize}

وزن نهایی هر کلمه در هر سند، حاصل‌ضرب این دو مقدار است:
\[
\text{TF-IDF}(t, d) = \text{TF}(t, d) \times \text{IDF}(t)
\]
این بازنمایی، پرکاربردترین روش کلاسیک برای تبدیل متن به بردار عددی است
و پایهٔ بسیاری از موتورهای جستجو و سیستم‌های بازیابی اطلاعات
نیز همین ایده است\footnote{\LR{Information Retrieval}}.

\section{کاربردهای مقدماتی}

پس از بازنمایی متن به‌صورت بردار عددی (با \LR{TF-IDF} یا کیسهٔ کلمات)،
می‌توان مستقیماً از الگوریتم‌های فصل‌های قبل استفاده کرد:

\begin{itemize}
    \item \textbf{طبقه‌بندی متن}: مثلاً تشخیص اسپم یا تحلیل احساسات
          (مثبت/منفی/خنثی بودن یک نظر مشتری) با استفاده از بیز ساده یا
          ماشین بردار پشتیبان.
    \item \textbf{خوشه‌بندی متن}: گروه‌بندی خودکار مقالات خبری بر اساس
          موضوع، بدون داشتن برچسب از پیش، با استفاده از k-means روی
          بردارهای \LR{TF-IDF}.
\end{itemize}

نکتهٔ مهم این است که بردارهای حاصل از متن معمولاً \textbf{ابعاد بسیار
بالا} (به تعداد کل کلمات منحصربه‌فرد واژگان) و \textbf{بسیار پراکنده}
دارند (اکثر درایه‌ها صفرند، چون هر سند فقط شامل بخش کوچکی از کل واژگان
است)؛ به همین دلیل، همان‌طور که در فصل \ref{ch:data} دیدیم، شباهت
کسینوسی معمولاً معیار مناسب‌تری نسبت به فاصلهٔ اقلیدسی برای مقایسهٔ این
بردارهاست.

\section{مقدمه‌ای بر داده‌کاوی وب}

\dmterm{\textbf{داده‌کاوی وب}}{Web Mining} به کاربرد تکنیک‌های
داده‌کاوی بر روی داده‌های وب اشاره دارد و معمولاً به سه زیرشاخه تقسیم
می‌شود:

\begin{itemize}
    \item \dmterm{\textbf{کاوش محتوای وب}}{Web Content Mining}: استخراج
          دانش از محتوای صفحات وب (متن، تصویر)، که عملاً تعمیمی از
          داده‌کاوی متن به مقیاس وب است.
    \item \dmterm{\textbf{کاوش ساختار وب}}{Web Structure Mining}: تحلیل
          گراف پیوندهای میان صفحات وب (مثلاً الگوریتم‌هایی شبیه
          \LR{PageRank} که برای رتبه‌بندی صفحات از ساختار پیوندها
          استفاده می‌کنند).
    \item \dmterm{\textbf{کاوش استفاده از وب}}{Web Usage Mining}: تحلیل
          رفتار کاربران در یک وب‌سایت (مانند مسیر کلیک‌ها) برای بهبود
          طراحی سایت یا سیستم‌های توصیه‌گر.
\end{itemize}

پرداختن کامل به این حوزه از حوصلهٔ این کتاب مقدماتی خارج است، اما آشنایی
با این تقسیم‌بندی برای شناخت جایگاه داده‌کاوی متن در تصویر بزرگ‌تر
مفید است.

\section*{پیاده‌سازی در پایتون}

پیاده‌سازی \texttt{TfidfVectorizer} و \texttt{CountVectorizer} از
\texttt{scikit-learn}، به‌همراه یک مثال کامل طبقه‌بندی متن، در بخش
«فصل ۹» پیوست الف (\ref{app:python}) آمده است.

\section*{خلاصهٔ فصل}

در این فصل با چالش اصلی کار با داده‌های متنی — نیاز به تبدیل متن به
بازنمایی عددی — آشنا شدیم. گام‌های پیش‌پردازش متن، کیسهٔ کلمات و
\LR{TF-IDF} را بررسی کردیم و دیدیم چگونه پس از این تبدیل، می‌توان از
همان الگوریتم‌های طبقه‌بندی و خوشه‌بندی فصل‌های قبل روی متن استفاده
کرد. در پایان، نگاهی کلی به سه زیرشاخهٔ داده‌کاوی وب انداختیم. در فصل
پایانی کتاب، به‌طور مقدماتی با داده‌های پیچیده‌تر (گراف و سری زمانی)
آشنا خواهیم شد.

\section*{تمرین‌ها}

\begin{enumerate}
    \item چرا نمی‌توان متن خام را مستقیماً به یک الگوریتم مانند درخت
          تصمیم یا k-means داد؟
    \item دو نمونه از ایست‌واژه‌های فارسی را نام ببرید و توضیح دهید چرا
          حذف آن‌ها معمولاً به بهبود کیفیت مدل کمک می‌کند.
    \item فرض کنید کلمهٔ «کامپیوتر» در یک سند خاص ۵ بار تکرار شده، اما
          در تقریباً همهٔ اسناد مجموعه نیز حضور دارد. وزن \LR{TF-IDF}
          این کلمه چه تفاوتی با کلمه‌ای کمیاب‌تر اما با همان فراوانی
          خواهد داشت؟
    \item چرا شباهت کسینوسی برای مقایسهٔ بردارهای \LR{TF-IDF} مناسب‌تر
          از فاصلهٔ اقلیدسی است؟
    \item تفاوت کاوش محتوای وب و کاوش ساختار وب را با یک مثال از هرکدام
          توضیح دهید.
\end{enumerate}

\chapter{مقدمه‌ای بر داده‌های پیچیده}
\label{ch:complex-data}

\section*{اهداف این فصل}
\begin{itemize}
    \item آشنایی مقدماتی با داده‌کاوی گراف
    \item آشنایی مقدماتی با تحلیل سری‌های زمانی
    \item آشنایی با مسیرهای پیشرفته‌تر برای ادامهٔ یادگیری
\end{itemize}

\section{فراتر از داده‌های جدولی}

در تمام فصل‌های پیشین، فرض بر این بود که داده به‌صورت جدولی است: هر
سطر یک نمونهٔ مستقل و هر ستون یک ویژگی. بسیاری از داده‌های واقعی، اما،
ساختاری فراتر از این دارند — روابط میان نمونه‌ها (مانند شبکه‌های
اجتماعی) یا ترتیب زمانی میان مشاهدات (مانند قیمت سهام). این فصل پایانی،
به‌صورت مقدماتی و صرفاً برای آشنایی، دو نوع از این داده‌های پیچیده‌تر
را معرفی می‌کند تا خواننده افق‌های بعدی مطالعه را بشناسد.

\section{داده‌کاوی گراف}

در بسیاری از مسائل، داده به‌طور طبیعی به‌صورت
\dmterm{\textbf{گراف}}{Graph} — مجموعه‌ای از
\dmterm{گره‌ها}{Nodes / Vertices} و
\dmterm{یال‌ها}{Edges} میان آن‌ها — بازنمایی می‌شود: کاربران یک شبکهٔ
اجتماعی و روابط دوستی میانشان، صفحات وب و پیوندهای میانشان، یا مولکول‌ها
و پیوندهای شیمیایی. \dmterm{\textbf{داده‌کاوی گراف}}{Graph Mining} به
کشف الگو در این نوع داده می‌پردازد.

برخی از سؤالات رایج در این حوزه عبارت‌اند از:

\begin{itemize}
    \item \dmterm{\textbf{تشخیص اجتماع}}{Community Detection}: یافتن گروه‌هایی از گره‌ها
          که درون خود پیوندهای متراکم‌تری نسبت به بیرون دارند — مفهومی
          که شباهت زیادی به خوشه‌بندی (فصل \ref{ch:clustering}) دارد،
          با این تفاوت که ورودی آن روابط میان نمونه‌هاست، نه ویژگی‌های
          عددی هر نمونه.
    \item \dmterm{\textbf{مرکزیت گره}}{Centrality}: سنجش میزان اهمیت یا نفوذ هر گره در
          گراف (مثلاً شناسایی کاربران تأثیرگذار در یک شبکهٔ اجتماعی).
    \item \dmterm{\textbf{پیش‌بینی پیوند}}{Link Prediction}: پیش‌بینی این‌که کدام یال‌های
          جدید در آینده به گراف اضافه خواهند شد (مثلاً پیشنهاد دوستی در
          شبکه‌های اجتماعی).
\end{itemize}

\section{تحلیل سری‌های زمانی}

\dmterm{\textbf{سری زمانی}}{Time Series} دنباله‌ای از مشاهدات است که
به ترتیب زمانی ثبت شده‌اند (مثلاً دمای روزانه، قیمت سهام، یا میزان
فروش ماهانه). تفاوت اساسی این نوع داده با داده‌های جدولی معمولی، وجود
\textbf{وابستگی زمانی} میان مشاهدات متوالی است؛ برخلاف فرض رایج در
بسیاری از الگوریتم‌های پیشین این کتاب که نمونه‌ها را مستقل از هم در
نظر می‌گرفتند.

\dmterm{\textbf{تجزیهٔ سری زمانی}}{Time Series Decomposition} یک سری
زمانی را معمولاً به سه مؤلفهٔ اصلی تجزیه می‌کند:

\begin{itemize}
    \item \dmterm{\textbf{روند}}{Trend}: تغییر کلی و بلندمدت (صعودی یا
          نزولی) در داده.
    \item \dmterm{\textbf{فصلی بودن}}{Seasonality}: الگوهای تکرارشونده
          در بازه‌های زمانی ثابت (مثلاً افزایش فروش هر سال در نزدیکی
          یک مناسبت خاص).
    \item \dmterm{\textbf{باقیمانده}}{Residual}: نوسانات تصادفی باقی‌مانده
          پس از کنار گذاشتن روند و فصلی بودن.
\end{itemize}

برای پیش‌بینی مقادیر آیندهٔ یک سری زمانی، از مدل‌های کلاسیکی مانند
\dmterm{\LR{ARIMA}}{AutoRegressive Integrated Moving Average} تا
روش‌های مدرن مبتنی بر یادگیری عمیق (مانند شبکه‌های بازگشتی\footnote{\LR{Recurrent Neural Networks (RNN)}})
استفاده می‌شود که بررسی
کامل آن‌ها خارج از حوصلهٔ این کتاب مقدماتی است.

\section{جمع‌بندی و مسیرهای پیشرفته‌تر}

در طول این کتاب، مسیری از مفاهیم پایه (فصل ۱ و ۲)، از میان یادگیری با
نظارت (طبقه‌بندی و رگرسیون در فصل‌های ۳ تا ۵)، تا یادگیری بدون نظارت
(قوانین انجمنی، خوشه‌بندی و تشخیص داده پرت در فصل‌های ۶ تا ۸)، و در
پایان مقدمه‌ای بر داده‌های متنی و پیچیده (فصل‌های ۹ و ۱۰) طی کردیم. این
کتاب صرفاً نقطهٔ شروع است؛ برای ادامهٔ مسیر، خواننده می‌تواند به سراغ
موضوعاتی مانند یادگیری عمیق، یادگیری تقویتی، پردازش زبان طبیعی مدرن (بر
پایهٔ ترنسفورمرها)، و مقیاس‌پذیری داده‌کاوی روی داده‌های عظیم (با
ابزارهایی مانند Spark) برود. مراجع اصلی این کتاب — Han/Kamber/Pei و
Tan/Steinbach/Kumar — نیز برای مطالعهٔ عمیق‌تر هر یک از فصل‌های این
کتاب به‌شدت توصیه می‌شوند.

\section*{تمرین‌ها}

\begin{enumerate}
    \item شباهت و تفاوت «تشخیص اجتماع» در گراف با «خوشه‌بندی» در
          فصل \ref{ch:clustering} را توضیح دهید.
    \item یک مثال واقعی از پیش‌بینی پیوند در یک شبکهٔ اجتماعی که با آن
          آشنایی دارید، بیان کنید.
    \item برای فروش ماهانهٔ یک فروشگاه لباس، مؤلفه‌های روند، فصلی بودن
          و باقیمانده را با یک مثال فرضی توضیح دهید.
    \item چرا فرض استقلال نمونه‌ها که در بسیاری از الگوریتم‌های فصل‌های
          قبل وجود داشت، برای داده‌های سری زمانی مناسب نیست؟
    \item با توجه به علایق خودتان، یکی از مسیرهای پیشرفته‌تر معرفی‌شده
          در این فصل (یادگیری عمیق، یادگیری تقویتی، پردازش زبان طبیعی
          مدرن، یا داده‌کاوی مقیاس‌بزرگ) را انتخاب کنید و توضیح دهید
          چرا برای ادامهٔ مطالعه به آن علاقه‌مندید.
\end{enumerate}


% ---------- پیوست‌ها ----------
\appendix
\chapter{پایتون برای داده‌کاوی}
\label{app:python}

این پیوست شامل پیاده‌سازی عملی مثال‌های هر فصل با استفاده از کتابخانه‌های \texttt{pandas}،
\texttt{numpy}، \texttt{scikit-learn}، \texttt{matplotlib} و \texttt{mlxtend} است.
کد کامل هر بخش در پوشهٔ texttt{code} مخزن گیت‌هاب قرار دارد و در اینجا تنها بخش‌های
کلیدی برای تطبیق با متن کتاب آورده شده‌اند.

\section{آماده‌سازی محیط}
\begin{lstlisting}
pip install pandas numpy scikit-learn matplotlib mlxtend
\end{lstlisting}

\section{فصل ۲ — پیش‌پردازش داده}
\begin{lstlisting}
import pandas as pd
from sklearn.preprocessing import StandardScaler
from sklearn.decomposition import PCA

df = pd.read_csv("data.csv")
df = df.dropna()

scaler = StandardScaler()
X_scaled = scaler.fit_transform(df.select_dtypes(include="number"))

pca = PCA(n_components=2)
X_pca = pca.fit_transform(X_scaled)
\end{lstlisting}

\section{فصل ۳ و ۴ — طبقه‌بندی}
\begin{lstlisting}
from sklearn.tree import DecisionTreeClassifier
from sklearn.naive_bayes import GaussianNB
from sklearn.neighbors import KNeighborsClassifier
from sklearn.svm import SVC
from sklearn.ensemble import RandomForestClassifier
from sklearn.model_selection import train_test_split
from sklearn.metrics import classification_report, roc_auc_score

X_train, X_test, y_train, y_test = train_test_split(X, y, test_size=0.2)

clf = DecisionTreeClassifier(max_depth=5)
clf.fit(X_train, y_train)
print(classification_report(y_test, clf.predict(X_test)))
\end{lstlisting}

\section{فصل ۶ — قوانین انجمنی}
\begin{lstlisting}
from mlxtend.frequent_patterns import apriori, association_rules

frequent_itemsets = apriori(basket_df, min_support=0.05, use_colnames=True)
rules = association_rules(frequent_itemsets, metric="confidence", min_threshold=0.6)
\end{lstlisting}

\section{فصل ۷ — خوشه‌بندی}
\begin{lstlisting}
from sklearn.cluster import KMeans, DBSCAN, AgglomerativeClustering

kmeans = KMeans(n_clusters=3, random_state=42).fit(X_scaled)
dbscan = DBSCAN(eps=0.5, min_samples=5).fit(X_scaled)
\end{lstlisting}

\section{فصل ۸ — تشخیص داده‌های پرت}
\begin{lstlisting}
from sklearn.ensemble import IsolationForest
from sklearn.neighbors import LocalOutlierFactor

iso = IsolationForest(contamination=0.05).fit(X_scaled)
lof = LocalOutlierFactor(n_neighbors=20)
outliers = lof.fit_predict(X_scaled)
\end{lstlisting}

\section{فصل ۹ — داده‌کاوی متن}
\begin{lstlisting}
from sklearn.feature_extraction.text import TfidfVectorizer

vectorizer = TfidfVectorizer(stop_words="english")
X_text = vectorizer.fit_transform(corpus)
\end{lstlisting}

\chapter{دیتاست‌های تمرینی}
\label{app:datasets}

در این پیوست فهرستی از دیتاست‌های عمومی و رایگان معرفی می‌شود که برای تمرین‌های هر فصل
مناسب هستند (مثلاً از UCI Machine Learning Repository و Kaggle). لینک دقیق هر دیتاست
و اسکریپت دانلود آن در پوشهٔ \texttt{data/} مخزن گیت‌هاب قرار می‌گیرد.

\begin{itemize}
    \item فصل ۳ و ۴ (طبقه‌بندی): Iris، Adult (Census Income)
    \item فصل ۵ (رگرسیون): Boston/California Housing
    \item فصل ۶ (قوانین انجمنی): Online Retail Dataset
    \item فصل ۷ و ۸ (خوشه‌بندی و داده پرت): Mall Customers، Credit Card Fraud
    \item فصل ۹ (متن): 20 Newsgroups
\end{itemize}


\backmatter
\printbibliography[title={مراجع}]

\end{document}
